\section{First-principles completion of the formal-local theorem}
\label{app:completion-formal-local-theorem}

The topology is part of the formal-Darboux theorem.  The residue dual
\(\mathfrak h^\vee_{\mathrm{cont}}\), the diagonal Casimir, the
Hamiltonian coadjoint action, the Hochschild equivariance homotopies, the
base heat-kernel homotopy, and any BV propagator datum live in one
completed filtered tensor category; outside that category the same
formulas are only polynomial identities.  The symbol \(\mathsf C\) below is
the coefficient--analytic presentation of the admissible host datum
\(\mathcal B^{\mathrm{adm}}\) of
Definition~\ref{defn:admissible-host-category}.  The theorem-level
admissible datum is \(\mathcal B^{\mathrm{adm}}\); \(\mathsf C\) records the
filtered nuclear Tate topology in which the local Hamiltonian BF theory,
the brane trace \(J\), and the Chevalley--Eilenberg/polyvector
dictionary are continuous.

The obstruction classes of the local theorem are then honest
cohomology classes in filtered \(L_\infty\) deformation complexes.
The algebraic classes govern separate lifts: radial Weyl compatibility,
Bochner--Martinelli current transfer, Tate bar-cobar admissibility,
Hochschild equivariance and centre-lift, the scalar-contact chain map,
quantum master equation counterterms, and the coefficient coevaluation
and brane-defect propagator datum.  A state extension adds the physical
\(N\to\infty\) trace-state row.  A theory \(\mathcal T\) compares its
anomalies with the local theorem by a morphism to the corresponding
completed deformation complex.  The algebraic comparison datum is
\(\mathfrak U_{\mathcal T}\); the local trace-sector Maurer--Cartan
element is \(\Theta_{\mathrm{loc}}\).

\subsection{The local presentation of the admissible host category}
\label{ssec:completion-host-category}

Tate algebra gives formal power series, residue duals, and infinite
Matlis sums.  Renormalisation gives heat kernels, wavefront
conditions, and distributional products against test functions.  The
common coefficient presentation is the filtered nuclear Tate dg
category.  It becomes an object of
\(\mathcal B^{\mathrm{adm}}\) only after the operation list, kernel list,
scalar convention, and Roos/stable-trace convention of
Definition~\ref{defn:admissible-host-category} are supplied.

\begin{defn}[Filtered nuclear Tate presentation]
\label{def:completion-host-category}
\(\mathsf C=\mathrm{Ch}^{\otimes,\mathrm{fil}}_{\mathrm{Tate,nuc},\C}\)
is the symmetric monoidal dg category presenting the
coefficient--analytic part of \(\mathcal B^{\mathrm{adm}}\).  Its
objects are complete decreasingly
filtered cochain complexes
\(V=(V,F^\bullet V,d_V)\) with
\[
  V\xrightarrow{\sim}\varprojlim_r V/F^r V,
  \qquad
  \bigcap_r F^r V=0,
\]
each associated graded
\(\mathrm{gr}^r_F V=F^r V/F^{r+1}V\) carried by a nuclear locally
convex cochain complex tensored with an admissible Tate
\(\C\)-coefficient module
(a linearly compact profinite-dimensional module, i.e. a strict
projective system of finite-dimensional \(\C\)-vector spaces, in the
sense of the local dictionary).  The completed tensor product is the strict
inverse limit of the term-wise nuclear projective tensor product on
the analytic factor with the term-wise Tate tensor product on the
coefficient factor:
\[
  V\widehat\otimes W
  =\varprojlim_r\bigl((V/F^rV)\otimes_\pi(W/F^rW)\bigr).
\]
The strict continuous dual is
\(V^\vee_{\mathrm{cont}}=\Hom^{\mathrm{cts}}_\C(V,\C)\), computed in
the same projective topology.  A BV kernel may instead use a declared
coefficient dual \(D_V\) with its own completed tensor product; when
\(D_V\ne V^\vee_{\mathrm{cont}}\), the statement concerns the declared
kernel pair and not the strict continuous dual.
\end{defn}

\begin{lem}[Coefficient axioms of \(\mathsf C\)]
\label{lem:completion-host-axioms}
\(\mathsf C\) satisfies the following four properties.  These are the
coefficient and analytic parts of
Definition~\ref{defn:admissible-host-category}(A1)--(A5), (A8), and
(A9); they do not include scalar reduction, stable trace passage, or the
theorem-specific operation and kernel lists.
\begin{description}
\item[\textnormal{C1 (filtered Tate exactness)}] For every inverse
  system \(\{V_M\}_{M\ge 0}\) used in the constructions of
  \S\ref{ssec:completion-promatlis}--\S\ref{ssec:completion-qme-tower},
  the transition maps are continuous and strict.  When ordinary
  inverse limits are used, Mittag--Leffler holds; when derived inverse
  limits \(\mathbf R\!\varprojlim\) are used, the corresponding
  \(\varprojlim^1\)-defects are recorded as obstruction classes in
  the master complex.  In every case the basic exact sequence
  \[
    0\to\varprojlim_M{}^1 H^{k-1}(V_M)\to
    H^k(\mathbf R\!\varprojlim_M V_M)\to
    \varprojlim_M H^k(V_M)\to 0
  \]
  holds.
\item[\textnormal{C2 (continuous and declared duality)}] For an
  admissible perfect pair \((V,V^\vee_{\mathrm{cont}})\) in
  \(\mathsf C\), the evaluation
  \(V^\vee_{\mathrm{cont}}\widehat\otimes V\to\C\) is continuous.  The
  coevaluation used by a BV propagator is a map
  \[
    \C\longrightarrow V\widehat\otimes D_V
  \]
  for the declared coefficient dual \(D_V\).  In the strict Tate pair
  \(D_V=V^\vee_{\mathrm{cont}}\).  In the weighted or
  kernel-admissible completions used for the diagonal Casimir,
  \(D_V\) may be a product coefficient module replacing the strict
  direct-sum dual.  Perfect evaluation and convergent Casimir are
  distinguished; the second map belongs to the declared completed
  category.
\item[\textnormal{C3 (local-functional subcategory)}]
  A subspace
  \(\mathcal O_{\mathrm{loc}}(\mathcal E)\subset\mathcal O(\mathcal E)\)
  of local functionals on a BV field complex is closed under \(Q\),
  the BV bracket \(\{-,-\}_{\mathrm{BV}}\), the regularized BV
  Laplacian \(\Delta_L\), Costello renormalization-group flow, and
  counterterm subtraction.  For a brane-defect theory the
  brane-local subspace is
  \(\mathcal O^{\mathrm{bd}}_{\mathrm{loc}}(\mathcal E_X;\mathcal E_L)\)
  consisting of functionals supported on \(X\), on \(L\), or on
  controlled collision strata involving \(L\).
\item[\textnormal{C4 (distributional current admissibility)}]
  The reduced principal-part current sector uses
  \(\mathcal D'_c(I)\widehat\otimes\mathcal P\) with multiplication
  \(C^\infty_c(I)\times\mathcal D'_c(I)\to\mathcal D'_c(I)\),
  \((a,B)\mapsto aB\), while the product
  \(\mathcal D'_c(I)\times\mathcal D'_c(I)\to\mathcal D'_c(I)\) is
	  excluded from the admissible operations.  The reduced principal-part bracket
	  of the brane current sector
	  factors through this restricted multiplication; the unreduced
  brane-defect QME requires a wavefront condition
  (\S\ref{ssec:completion-graph-qme}).
\end{description}
\end{lem}

\begin{proof}
C1 is the Milnor sequence
\cite{milnor-axiomatic-homology}.  C2 is the distinction used here
between a continuous nondegenerate evaluation pairing
\(E\widehat\otimes E^\vee\to\C\) in linearly compact / Tate spaces and
the coevaluation tensor \(\C\to E\widehat\otimes E^\vee\) needed to
write the inverse BV kernel: a continuous evaluation map exists for any
continuous dual, but the coevaluation map exists only after the
convergent regulator of
\S\ref{ssec:completion-bracket-kernel-admissible}.  C3 is the local
functional subcategory of Costello \cite{costello-renormalization}*{Ch.~5}
adapted to brane-defect support strata; closure under \(\Delta_L\) is
the Costello effective-action calculus and closure under counterterm
subtraction follows from the regulator-independence of local
divergences.  C4 is H\"ormander's product theorem
\cite{hormander-vol1}*{Thm.~8.2.10}: the smooth-times-distribution
product is everywhere defined, while two distributions multiply when no
covectors \((x,\xi)\in\operatorname{WF}(u)\) and
\((x,-\xi)\in\operatorname{WF}(v)\) occur.  The reduced
	principal-part bracket of the brane current sector respects exactly this
	restriction.
\end{proof}

All subsequent analytic constructions are performed in \(\mathsf C\).
They are statements in \(\mathcal B^{\mathrm{adm}}\) only after the
relevant operation list \(\mathcal O_{\mathrm{adm}}\), kernel list
\(\mathcal K_{\mathrm{adm}}\), scalar convention, and Roos convention are
attached.

\subsection{The local Hamiltonian BF theory as a genuine BV theory}
\label{ssec:completion-bv-theorem}

The classical master equation \(\{S_H,S_H\}_{\mathrm{BV}}=0\) of the
formal-local Hamiltonian BF theory holds in the completed
local-functional space as an honest local-functional identity.  The proof
uses the classical BV bracket and avoids the Costello renormalization complex and the
brane-defect propagator; it is the foundation on which every quantum
construction below depends.

\begin{thm}[Classical master equation]
\label{thm:completion-cme}
Let \(A=\C[[z_1,z_2]]\), \(\mathfrak m=(z_1,z_2)\),
\(\mathfrak h=A/\C\cdot 1\), and let
\(\mathcal P=\mathfrak h^\vee_{\mathrm{cont}}=
\bigoplus_{a+b>0}\C\rho_{a,b}\) with
\(\rho_{a,b}=z_1^{-a-1}z_2^{-b-1}\,dz_1\,dz_2\).  Define
\[
  \mathcal E_H(U\times B)
  =\Omega^\bullet_c(U)\widehat\otimes
   \Omega^{0,\bullet}_c(B)\widehat\otimes
   \bigl(\mathfrak h[1]\oplus\mathcal P[2]\bigr)
\]
on \(U\subset\R^2_{t,s}\), \(B\subset\C^2_{z_1,z_2}\), with
differential \(Q_H=d_{\R^2}+\bar\partial_{\C^2}+d_{\mathfrak g}\),
\((-1)\)-shifted symplectic pairing
\[
  \omega_H(\alpha\otimes f,\beta\otimes\rho)
  =\int_{\R^2\times\C^2}\alpha\wedge\beta\wedge dz_1\,dz_2\,
   \cdot\operatorname{Res}_0(f\rho),
\]
and classical action
\[
  S_H(\mathcal A,\mathcal B)
  =\int_X\bigl\langle\mathcal B,
    (d+\bar\partial)\mathcal A
    +\tfrac12[\mathcal A,\mathcal A]_{\mathfrak h}\bigr\rangle.
\]
Then \(\{S_H,S_H\}_{\mathrm{BV}}=0\) in
\(\mathcal O^{\mathrm{bd}}_{\mathrm{loc}}(\mathcal E_H)\).
\end{thm}

\begin{proof}
The bracket expansion of \(\{S_H,S_H\}\) is the Maurer--Cartan
Bianchi identity for the connection \((d+\bar\partial)\mathcal A
+\tfrac12[\mathcal A,\mathcal A]_{\mathfrak h}\):
\[
  D\bigl(D\mathcal A+\tfrac12[\mathcal A,\mathcal A]\bigr)
  +[\mathcal A,D\mathcal A+\tfrac12[\mathcal A,\mathcal A]]
  =0.
\]
The four ingredients are \(D^2=0\)
(\((d+\bar\partial)^2=0\) on \(\R^2\times\C^2\)), the graded Leibniz
rule \(D[\mathcal A,\mathcal A]=2[D\mathcal A,\mathcal A]\), the
Jacobi identity in \(\mathfrak h\), and continuity of the residue
pairing.  Continuity is the C2 axiom: the Tate dual
\(\mathcal P\subset\mathfrak h^\vee_{\mathrm{cont}}\) is paired with
\(\mathfrak h\) by the continuous functional
\(\langle f,\rho\rangle=\operatorname{Res}_0(f\rho)\), which descends
through the integration by parts identity
\(\int_X\langle D\mathcal B,\mathcal A\rangle
+(-1)^{|\mathcal B|}\int_X\langle\mathcal B,D\mathcal A\rangle=0\)
for compactly supported test forms.  Therefore the bracket
\(\{S_H,S_H\}\) reduces to a sum of terms each of which vanishes by
one of the four ingredients.  The resulting equality holds in the
Tate-completed local functional space, since each ingredient is a
continuous operation in \(\mathsf C\).
\end{proof}

\begin{rmk}[BV degree, Koszul sign, propagator weight]
\label{rmk:completion-bv-conventions}
The conventions are those of \(d=\dim_\C\C^2=2\) holomorphic dimensions and
\(d_{\R^2}=2\) topological dimensions; BV degree
\(\deg_{\mathrm{BV}}\mathcal A=1\), \(\deg_{\mathrm{BV}}\mathcal B=2\)
following Costello's BV formalism
\cite{costello-renormalization}*{Chs.~2--4}; the propagator
\(P_{\varepsilon,L}\) has BV degree \(-1\); the holomorphic volume
form is \(dz_1\,dz_2\) (purely holomorphic); the worldline
\(L=\R_t\subset X\) is unframed in the local model and acquires the
\(S^3\)-framing only after a matched-conventions transfer
(\S\ref{app:comparison-conventions}).
\end{rmk}

\subsection{Hamiltonian shears and the trace uniqueness theorem}
\label{ssec:completion-trace-uniqueness}

The trace uniqueness theorem (Theorem~\ref{thm:Jf-coordinate-uniqueness})
states that \(J(z_1^az_2^b)=\Tr(\phi_1^a\phi_2^b)\) for every
bidegree.  The proof uses linear shears to relate same-degree
coefficients and one nonlinear Hamiltonian shear to relate different
total degrees.  Linear shears alone leave one constant per total
degree undetermined; the nonlinear shear closes the proof.

\begin{thm}[Hamiltonian-shear determination of coefficients]
\label{thm:completion-shear-determination}
Let
\[
  J^{\mathrm{st}}\colon\mathfrak h\longrightarrow
  A_{\partial,H}^{\mathrm{st}}
\]
be a continuous linear map to the stable scalar-reduced Hamiltonian
trace algebra.  Suppose that, in every Taylor-degree coefficient
window \(W\), its value is represented in ranks \(N\geq W\) by a
primitive \(GL_N\)-invariant single-trace class on the Dirac-derived
zero fibre, is natural under formal symplectomorphisms of
\((\C^2,dz_1\wedge dz_2)\), and satisfies
\[
  J^{\mathrm{st}}(z_i)
  =\overline{\Tr\phi_i}.
\]
Then
\[
  J^{\mathrm{st}}(z_1^az_2^b)
  =\overline{\Tr(\phi_1^a\phi_2^b)}
\]
in the degree-\((a+b)\) stable trace system for every \(a,b\geq0\)
with \(a+b>0\).  At fixed rank \(N\), the same formula defines the
class \(J_N^s(z_1^az_2^b)\) and retains the finite-\(N\) trace
relations and scalar-contact cocycle.
\end{thm}

\begin{proof}
Fix \(d=a+b\) and choose a representing rank \(N\geq d\).
Procesi's first fundamental theorem \cite{procesi}*{Thm.~1.3},
primitive block-additivity, and the Koszul homotopy for adjacent
transpositions on the moment-map zero fibre
(Lemma~\ref{lem:koszul-adjacent-transpositions}) give
\[
  J^{\mathrm{st}}(z_1^az_2^b)
  =c_{a,b}\overline{\Tr(\phi_1^a\phi_2^b)}
\]
for some scalar \(c_{a,b}\in\C\).  It suffices to prove
\(c_{a,b}=1\) for every \((a,b)\) with \(a+b>0\); the linear case
\(c_{1,0}=c_{0,1}=1\) is the normalization hypothesis
\(J^{\mathrm{st}}(z_i)=\overline{\Tr\phi_i}\).

\emph{Linear shears within fixed total degree.}  The Darboux flow
generated by \(H_{\mathrm{lin}}=z_1^2/2\), at time \(t\), is
\(z_1\mapsto z_1\), \(z_2\mapsto z_2+tz_1\), so
\((\Phi^*z_1^az_2^b)=\sum_{j=0}^b\binom{b}{j}t^jz_1^{a+j}z_2^{b-j}\).
Equivariance
\(J^{\mathrm{st}}(\Phi^*f)=\Phi^*J^{\mathrm{st}}(f)\) on the
stable brane trace system gives
\(\Phi^*\overline{\Tr(\phi_1^a\phi_2^b)}=\sum_{j}\binom{b}{j}t^j
\overline{\Tr(\phi_1^{a+j}\phi_2^{b-j})}\), and comparing
coefficients of \(t^j\) yields
\(c_{a+j,b-j}=c_{a,b}\) within fixed total degree \(a+b\).

\emph{Nonlinear shear across degrees.}  Consider the Hamiltonian
\(H_k=z_2^{k+1}/(k+1)\), whose flow is
\(z_1\mapsto z_1-tz_2^k\), \(z_2\mapsto z_2\).  Apply
Darboux-naturality to the linear observable \(z_1\), already
normalized to \(c_{1,0}=1\):
\[
  J^{\mathrm{st}}(z_1-tz_2^k)
  =\overline{\Tr(\phi_1)}
  -t\,c_{0,k}\overline{\Tr(\phi_2^k)} ,
\]
while the brane-side flow gives
\(\overline{\Tr(\phi_1)}-t\overline{\Tr(\phi_2^k)}\).  Comparing
coefficients of \(t\) yields \(c_{0,k}=1\), and linear shears
propagate this to \(c_{a,b}=1\) for every \((a,b)\) with \(a+b=k\).
Since \(k\ge 1\) is arbitrary the assertion follows for every total
degree.  Procesi \cite{procesi}*{Thm.~4.6} and Razmyslov
\cite{razmyslov}*{Thm.~2} govern the additional trace relations in
each fixed rank; the stable statement is their degreewise compatible
range.
\end{proof}

\begin{rmk}[Primitive trace relation]
\label{rmk:completion-multiplicativity-relation}
The hypothesis ``primitive single-trace'' is the
Loday--Quillen--Tsygan condition that the Lie algebra map \(J\) lands
in the connected primitive sector of the cyclic invariant algebra
\cite{loday-quillen}*{Thm.~6.2}\cite{tsygan}*{Thm.~1}; the
degreewise stable scalar-reduced use here is the manuscript's
separated-block adaptation.  It produces the trace word for each
monomial.  The linear normalization
\(c_{1,0}=c_{0,1}=1\) and the nonlinear shears determine its
coefficient in every total degree.  The resulting relation combines
the Procesi trace word, the Koszul resolution of
\([\phi_1,\phi_2]=0\), and single-trace primitivity.  The fixed-rank
classes \(J_N^s(f)\) remain in their finite-\(N\) trace quotients;
the map \(J^{\mathrm{st}}\) is their degreewise stable scalar-reduced
passage.
\end{rmk}

\subsection{\texorpdfstring{The decorated PBW--Stokes complex and \(\Omega^{\mathrm{rad}}_{a,b}=0\)}{The decorated PBW-Stokes complex and Omega-rad-a,b = 0}}
\label{ssec:completion-decorated-pbw-stokes}

The all-bidegree Weyl/radial compatibility criterion
(Theorem~\ref{thm:phi-hbar-all-order}) depends on the
Statement~\ref{stmt:app-radial-parts-hypotheses} clauses (1)
and (3).  Clauses (1)--(2) are the cited Harish-Chandra, Wallach, and
Levasseur--Stafford radial homomorphism, regular-Cartan landing, and
kernel input.  Clause (3) is the additional exact Weyl/Capelli
trace-image datum
\(\operatorname{Rad}_0(J_N(f))=\Delta^{-1}S_N(f)\Delta\) in the
normalization of the formal-Darboux theorem.  The decorated PBW--Stokes
complex below, in which the per-bidegree obstruction
\(\Omega^{\mathrm{rad}}_{a,b}=[(T_{a,b},E^+_{a,b})]\in
\operatorname{coker}B_{a,b}\) is tested against an explicit
contracting homotopy, is the finite-rank obstruction complex.

\begin{defn}[Decorated PBW--Stokes complex]
\label{def:completion-pbw-stokes}
Fix \(a,b\ge 0\) with \(a+b>0\).  Let \(\mathcal W_N\) be the matrix
Weyl algebra on \(X^i_j,Y^k_\ell\) of
Definition~\ref{def:app-radial-matrix-weyl}, with quantum moment-map
ideal \(I_N=\mathcal W_N\,\mu(\mathfrak{sl}_N)\).  Let
\(C^0_{a,b}=\bigl(\Omega^{a+b+2}_{\mathrm{PBW}}(X,Y;\hbar_{\mathrm{W}})\bigr)
\oplus\C\), the rank-\(b+2\) free \(\mathcal W_N\)-module of PBW
\((Y^jX^aY^{b-j})\)-words decorated by Vandermonde--Stokes
generators
\(s^{(\alpha)}_{a,b}\in C^0_{a,b}\) carrying the boundary
contributions of the type-\(A\) Dunkl--Calogero
correction (Lemma~\ref{lem:app-radial-dunkl-calogero-cancellation}).
Let \(C^{-1}_{a,b}=\mathcal W_N^{\oplus(b+1)}\) be the corresponding
free module of homotopy primitives, and define
\[
  B_{a,b}:C^{-1}_{a,b}\longrightarrow C^0_{a,b},
  \qquad
  B_{a,b}(\xi^{(j)})_{j=0}^b
  =\Bigl(\sum_{j=0}^bA^{(j)}_{a,b,N}\xi^{(j)},\,
   \sum_{j,\alpha}\partial^{(\alpha)}_{j,a,b}\xi^{(j)}\Bigr),
\]
where \(A^{(j)}_{a,b,N}\) are the primitive coefficients in the
Vandermonde-conjugated radial linear system
(\S\ref{ssec:completion-decorated-pbw-stokes-primitives}) and
\(\partial^{(\alpha)}_{j,a,b}\) are the Stokes-decoration boundary
operators recording the
Vandermonde-pole-stripping data on the \(\alpha\)-th coordinate
chart of the regular Cartan.  The transported shear vector is
\((T_{a,b},E^+_{a,b})\in C^0_{a,b}\) with
\(T_{a,b}=\sum_j\Tr(Y^jX^aY^{b-j})\) and \(E^+_{a,b}\) the
\(J_N(\{z_1^{a+1}/(a+1),z_2^{b+1}\}_{\hbar_{\mathrm{W}}})\) shear-error term.
\end{defn}

\subsubsection{\texorpdfstring{Primitive linear systems \(A^{(j)}_{a,b,N}\) from the Vandermonde-conjugated radial map}{Primitive linear systems A-j-a,b,N from the Vandermonde-conjugated radial map}}
\label{ssec:completion-decorated-pbw-stokes-primitives}

Under the homomorphism and injectivity hypotheses of
Statement~\ref{stmt:app-radial-parts-hypotheses}(1)--(2), the
Vandermonde-conjugated radial map sends \(X^i_j,Y^k_\ell\) into
formal differential operators on the regular Cartan in the
Vandermonde gauge \(\nabla^\Delta_i\) of
Definition~\ref{def:app-radial-vandermonde-gauge-variables}.  In this
gauge the Calogero--Dunkl correction terms cancel pairwise
(Lemma~\ref{lem:app-radial-dunkl-calogero-cancellation}), so the
radial image of \(\sum_{j=0}^b\Tr(Y^jX^aY^{b-j})\) coincides with
\(\Delta^{-1}\sum_{i=1}^N\operatorname{Op}_{\mathrm W}(z_1^az_2^b)
(\lambda_i,-\hbar_{\mathrm{W}}\partial_{\lambda_i})\Delta\) up to a Stokes
correction supported on the boundary of the regular Cartan.

\begin{defn}[Stokes primitive datum in the Vandermonde gauge]
\label{lem:completion-stokes-correction}
Fix \(a,b\ge 0\) with \(a+b>0\).  A Stokes primitive datum in this
bidegree is a collection of differential operators
\(A^{(j)}_{a,b,N}\in\mathcal W_N\) together with a Stokes correction
\(\tau_{a,b}\in C^{-1}_{a,b}\) satisfying
\[
  E_{a,b,N}=\sum_{i,j}A^{(j)}_{a,b,N}\mu^j_i+B_{a,b}(\tau_{a,b}),
\]
where \(\mu^j_i\) are generators of the quantum moment-map ideal.
When such data exist, the bidegree \((a,b)\) is
PBW--Stokes-solvable.  The exact rational calculations of
Appendix~\ref{sec:app-radial-parts-moyal} construct these data in the
recorded finite ranges.  Outside those ranges, existence of the datum is
the bidegree obstruction problem.
\end{defn}

\begin{rmk}[Vandermonde-gauge linear system]
Under Statement~\ref{stmt:app-radial-parts-hypotheses}(1)--(2), the
Vandermonde gauge variables satisfy
\([\lambda_i,\nabla^\Delta_j]=\hbar_{\mathrm{W}}\delta_{ij}\).  Products of
\(\nabla^\Delta_i\) expand into \((-\hbar_{\mathrm{W}}\partial_{\lambda_i})\)-terms
plus Vandermonde boundary terms; the two-body terms cancel pairwise by
Lemma~\ref{lem:app-radial-dunkl-calogero-cancellation}.  In a fixed
bidegree this leaves a finite rational linear system in the unknowns
\((A^{(j)}_{a,b,N},\tau_{a,b})\).  A solution is exactly a Stokes
primitive datum of Definition~\ref{lem:completion-stokes-correction}.
The manuscript does not assert a uniform solution in every bidegree
from this gauge computation alone; the uniform solution is the
all-bidegree radial primitive problem.
\end{rmk}

\begin{thm}[Decorated radial reformulation of the all-bidegree problem]
\label{thm:completion-omega-rad-vanishes}
On the decorated PBW-Stokes complex
\(C^\bullet_{a,b}\) of Definition~\ref{def:completion-pbw-stokes},
the per-bidegree obstruction
\(\Omega^{\mathrm{rad}}_{a,b}=[(T_{a,b},E^+_{a,b})]\in
\operatorname{coker}B_{a,b}\) vanishes if and only if there exists a
Stokes primitive datum
\((A^{(j)}_{a,b,N},\tau_{a,b})\) in the sense of
Definition~\ref{lem:completion-stokes-correction}, equivalently if and
only if the radial image of the trace map satisfies
clause~(3) of Statement~\ref{stmt:app-radial-parts-hypotheses}.  In
either formulation the all-bidegree hypothesis is converted into a
finite-rank exact-rational existence problem at each bidegree
\((a,b)\), independent of the radial normalization beyond
clauses~(1)--(2).  The contracting-homotopy formulation
(Route~A of Remark~\ref{rmk:completion-dual-kernel}) and the dual
kernel formulation (Route~B) are equivalent, and the finite
exact-rational calculations of Appendix~\ref{sec:app-radial-parts-moyal}
verify both routes through the recorded ranges of \((a,b)\); the
all-\((a,b)\) statement is the uniform existence of these
primitives on the entire decorated bicomplex.
\end{thm}

\begin{proof}
By Definition~\ref{lem:completion-stokes-correction}, the existence of
\((A^{(j)}_{a,b,N},\tau_{a,b})\) gives
\(E_{a,b,N}-B_{a,b}(\tau_{a,b})=\sum_{i,j}A^{(j)}_{a,b,N}\mu^j_i\in
I_N\), so \([(T_{a,b},E^+_{a,b})]=0\) in
\(\operatorname{coker}B_{a,b}\).  Conversely, if the class
vanishes, the cokernel computation produces such primitives by
finite-rank linear algebra in the decorated bicomplex.  The
correspondence with clause~(3) of
Statement~\ref{stmt:app-radial-parts-hypotheses} is the equality
\(\operatorname{Rad}_0(J_N(f))=\Delta^{-1}S_N(f)\Delta\) read off
from the Vandermonde-gauge expansion: the symmetric sum
\(\sum_i\operatorname{Op}_{\mathrm W}(f)(\lambda_i,\nabla^\Delta_i)\)
equals \(\Delta^{-1}S_N(f)\Delta\) by
Lemma~\ref{lem:app-radial-dunkl-calogero-cancellation}, so the
radial image of \(J_N(f)\) is determined by clause~(3) up to the
quantum moment-map ideal, hence by injectivity (clause~(2)).
Equivalence of Routes~A and~B follows from the duality
\(\operatorname{coker}B_{a,b}=(\ker B^*_{a,b})^\vee\) on the
finite-rank decorated bicomplex.
\end{proof}

\begin{rmk}[All-bidegree obstruction]
\label{rmk:completion-omega-rad-status}
Theorem~\ref{thm:completion-omega-rad-vanishes} converts the
Statement~\ref{stmt:app-radial-parts-hypotheses}(3) into a
constructible primitive-existence problem on a finite-rank
decorated bicomplex.  The remaining datum is the uniform
construction of \((A^{(j)}_{a,b,N},\tau_{a,b})\) for all
\((a,b)\); the finite exact-rational calculations of
Appendix~\ref{sec:app-radial-parts-moyal} give these primitives in the
recorded ranges, and
the all-\((a,b)\) statement is the assertion that the bicomplex
has a Lagrangian splitting compatible with the shears.  This
reformulation is mathematically stronger than the radial
normalization because it identifies the obstruction with a
\(\operatorname{coker}\)-class on a concrete finite-rank Weyl
target and identifies the all-\(N\) construction problem as the
existence of a
contracting homotopy or a dual-kernel certification at each
bidegree.
\end{rmk}

\begin{rmk}[Route B: dual kernel criterion]
\label{rmk:completion-dual-kernel}
The contracting homotopy of
Theorem~\ref{thm:completion-omega-rad-vanishes} is Route A of the
two acceptable proof routes; the equivalent Route B is the dual
kernel criterion
\((\varphi,-\lambda)\in\ker B^*_{a,b}\Rightarrow
\lambda(E^+_{a,b})-\varphi(T_{a,b})=0\), proved by pairing both
sides with the same Vandermonde-conjugated functionals.  Either
route closes the radial obstruction in the bidegree under the
homomorphism and injectivity hypotheses.
\end{rmk}

\subsection{The pro-Matlis weighted K\"othe envelope}
\label{ssec:completion-promatlis}

The attempted strict all-window Bochner--Martinelli current transfer
meets an obstruction in every finite window, given by the
residue-coadjoint leak
\(z_1^{N+2}\cdot\rho_{N+1,0}=-(N+2)\rho_{0,1}\)
(Theorem~\ref{thm:bmk-finite-window-obstruction}).  The replacement target is
the pro-Matlis tower; its analytic content is a continuous
Hamiltonian action, which requires a topology stronger than the bare
inverse limit of the finite Matlis windows.

\begin{defn}[Weighted K\"othe Hamiltonian space]
\label{def:completion-weighted-kothe}
Fix \(q>1\).  The weighted K\"othe Hamiltonian space is
\[
  \mathfrak h_q
  =\Bigl\{f=\sum_{a+b>0}f_{a,b}z_1^az_2^b\in\C[[z_1,z_2]]\,:\,
   p_n(f)=\sup_{a,b}|f_{a,b}|q^{a+b}(1+a+b)^n<\infty
   \,\,\forall\,n\Bigr\},
\]
a Fr\'echet algebra under
\(\{p_n\}_{n\ge 0}\) and the Poisson bracket of \(\mathfrak h\).
The declared pro-Matlis coefficient dual is the weighted principal-part
space
\[
  \mathcal P^\Pi_q
  =
  \Bigl\{\rho=\sum_{a+b>0}r_{a,b}\rho_{a,b}\,:\,
  p^*_n(\rho)=
  \sup_{a,b}|r_{a,b}|\,q^{-(a+b)}(1+a+b)^n<\infty
  \,\,\forall\,n\Bigr\}.
\]
The residue pairing
\(\langle\rho_{a,b},z_1^mz_2^n\rangle=\delta_{a,m}\delta_{b,n}\)
is continuous after increasing the Hamiltonian seminorm index:
for every \(n\) there is \(N\) such that
\(|\langle\rho,f\rangle|\le C_{n,N}p^*_n(\rho)p_N(f)\).
Thus \(\mathcal P^\Pi_q\) is the declared coefficient dual used by the
pro-Matlis target.  The continuous coefficient maps
\(\pi_N\colon\mathcal P^\Pi_q\to\mathcal P_{\leq N}\) give its
finite Matlis windows; the weighted seminorms specify its inclusion in
the bare coefficient product.
\end{defn}

\begin{prop}[Continuity estimate]
\label{prop:completion-promatlis-continuity}
There exist constants \(C_n\ge 0\) and integers
\(r=r(n)\ge 0\) such that for all \(f\in\mathfrak h_q\) and
\(\rho\in\mathcal P^\Pi_q\),
\[
  p^*_n(f\cdot\rho)\le C_n\,p_{n+r}(f)\,p^*_{n+r}(\rho),
\]
where \(f\cdot\rho\) is the residue coadjoint action.  Consequently
the bracket
\(\mathfrak h_q\widehat\otimes\mathcal P^\Pi_q\to\mathcal P^\Pi_q\)
is continuous, the action descends to a continuous
\(\mathfrak h_q\)-module structure, and the finite-window coherence law
\[
  \pi_N(f_{\le s}\cdot\xi)
  =
  \pi_N(f_{\le s}\cdot\pi_{N+s}\xi)
\]
holds for every polynomial Hamiltonian \(f_{\le s}\) of total degree
\(\le s\).  For a general \(f\in\mathfrak h_q\), the left side is the
limit of these polynomial truncation actions in the pro-Matlis topology.
\end{prop}

\begin{proof}
The residue coadjoint action is
\[
  z_1^pz_2^\ell\cdot\rho_{a,b}
  =-(pb-\ell a+p-\ell)\rho_{a-p+1,b-\ell+1}
\]
with negative-index terms zero
(Theorem~\ref{thm:principal-part-coadjoint-uniqueness}).  Put
\(P=p+\ell\), \(D=m+n\), and
\(R=a+b=D+P-2\).  The coefficient of \(\rho_{m,n}\) in
\(f\cdot\rho\) is
\[
  (f\cdot\rho)_{m,n}
  =
  -\!\!\sum_{\substack{p,\ell\ge0\\p+\ell>0}}
  (p(n+\ell-1)-\ell(m+p-1)+p-\ell)
  f_{p,\ell}r_{m+p-1,n+\ell-1},
\]
where summands with negative indices are omitted.  The coefficient in
absolute value is \(|pn-\ell m|\), hence is bounded by
\((1+D)(1+P)\).  For \(r\ge4\),
\[
  |f_{p,\ell}|
  \le p_{n+r}(f)\,q^{-P}(1+P)^{-(n+r)}
\]
and
\[
  |r_{m+p-1,n+\ell-1}|
  \le p^*_{n+r}(\rho)\,q^{R}(1+R)^{-(n+r)} .
\]
Multiplication by the output seminorm weight gives the exponential
factor
\[
  q^{-D}q^{-P}q^R=q^{-2},
\]
because \(R=D+P-2\).  Hence
\[
\begin{aligned}
 & |(f\cdot\rho)_{m,n}|q^{-D}(1+D)^n  \\
 &\quad \le
 C\,p_{n+r}(f)p^*_{n+r}(\rho)
 (1+D)^n
 \sum_{P\ge1}
 (P+1)(1+D)(1+P)
 (1+P)^{-(n+r)}(1+D+P)^{-(n+r)} .
\end{aligned}
\]
The displayed sum is bounded uniformly in \(D\).  Indeed, for
\(P\le D+1\) it is bounded by a constant times
\((1+D)^{1-r}\sum_{P\ge1}(P+1)(1+P)^{1-(n+r)}\), and for
\(P>D+1\) it is bounded by a constant times
\((1+D)^{n+1}\sum_{P>D}P^{2-2(n+r)}\), which is
\(O((1+D)^{4-n-2r})\).
For \(r\ge4\) both bounds are uniform in \(D\).  Taking the supremum
over \(m,n\) gives the asserted seminorm estimate.

The estimate is bilinear continuity of the residue coadjoint action in
the declared weighted pro-Matlis topology.  The finite-window coherence
law follows because \(\pi_N(f_{\le s}\cdot\xi)\) depends only on those
\(\xi_{a,b}\) with \(a+b\le N+s\) for polynomial \(f_{\le s}\).  For
general \(f\in\mathfrak h_q\), polynomial truncations converge to \(f\)
in every \(p_n\), and the seminorm estimate carries this convergence to
the action on \(\mathcal P^\Pi_q\).
\end{proof}

\begin{thm}[Native pro-Matlis replacement target]
\label{thm:completion-native-promatlis}
The pair
\((\mathfrak h_q,\mathcal P^\Pi_q)\) is the chosen replacement target
for formal Hamiltonian actions on residue-current systems with
Bochner--Martinelli leakage.  It carries a continuous coadjoint action
and the finite-window coherence law of
Proposition~\ref{prop:completion-promatlis-continuity}.  The category
\(\mathrm{Curr}_{\mathrm{BM}}\) used here consists of inverse systems
of coefficient spaces
\((\rho_N:\mathcal P_N\twoheadrightarrow\mathcal P_{N-1})\)
carrying this exact window-loss identity and a continuous full action
on the inverse-limit object.  A terminality statement is a separate
universal-property theorem in this category.
\end{thm}

\begin{proof}
The truncations \(\pi_N:\mathcal P^\Pi_q\to\mathcal P_{\le N}\) satisfy
the finite-window coherence law
\(\pi_N(f_{\le s}\cdot\xi)=\pi_N(f_{\le s}\cdot\pi_{N+s}\xi)\) of
Proposition~\ref{prop:completion-promatlis-continuity}.  Continuity of
the bracket on \(\mathcal P^\Pi_q\) gives a coherent action on the
weighted pro-tower.  The maps \(\pi_N\) are coefficient projections
with exact polynomial window loss; the full module structure resides
on \(\mathcal P^\Pi_q\), where formal Hamiltonians act as limits of
their polynomial truncations.
\end{proof}

\begin{rmk}[The category determines the replacement datum]
\label{rmk:completion-currbm-category}
The replacement datum consists of
\(\mathcal P^\Pi_q\), its coefficient projections, and its
window-loss coherence law.  Direct sums
\(\mathcal P^\Pi_q\oplus W\) with spectator modules distinguish this
specified datum from a terminal object of
\(\mathrm{Curr}_{\mathrm{BM}}\).
\end{rmk}

\subsection{\texorpdfstring{Bracket-admissible \(B\) and kernel-admissible \(\mathcal K_B\)}{Bracket-admissible B and kernel-admissible K-B}}
\label{ssec:completion-bracket-kernel-admissible}

Levels 2 and 3 of the CE/PV ladder hold under a
bracket-admissible subspace
\(B\subset\mathrm{PV}_{\mathrm{coord}}(A^{\mathrm{cl}}_{\partial,\mathrm{coord}})\)
and a kernel-admissible completion \(\mathcal K_B\) in which the
diagonal Casimir converges.  The pro-Matlis weighted
K\"othe topology gives both.  The symbol
\(\mathrm{PV}_{\mathrm{coord}}\) denotes algebraic coordinate
polyvectors on the completed trace coordinate ring, with cotangent
Schouten bracket of degree \(-1\); it is not the geometric Dolbeault
polyvector complex of the compact BCOV target.

Put \(P_q=\mathcal P^\Pi_q\) and
\(\mathfrak g_q=\mathfrak h_q\ltimes P_q[1]\).  The standard
Chevalley--Eilenberg cochain algebra has generators
\[
  c^\rho\quad(\rho\in P_q),\qquad
  u_f\quad(f\in\mathfrak h_q),\qquad
  |c^\rho|_{\mathrm{CE}}=1,\quad
  |u_f|_{\mathrm{CE}}=2.
\]
Give \(u_f\) cotangent weight one and \(c^\rho\) cotangent weight
zero.  The coordinate regrading
\[
  |a|_{\mathrm{coord}}
  =|a|_{\mathrm{CE}}-2\operatorname{wt}_u(a)
\]
places \(u_f\) in degree zero.  With
\(O_f=J^{\mathrm{st}}(f)\), the completed coordinate map is
\[
  \Phi_{\mathrm{coord}}\colon
  \operatorname{Regr}_u
  C^\bullet_{\mathrm{CE},\mathrm{coord}}(\mathfrak g_q)
  \longrightarrow
  \mathrm{PV}_{\mathrm{coord}}
    (A_{\partial,H}^{\mathrm{st}}),
  \qquad
  c^\rho\longmapsto\partial_\rho,\quad
  u_f\longmapsto O_f,
\]
where
\[
  \partial_\rho(O_f)=\langle\rho,f\rangle,
  \qquad
  \Phi_{\mathrm{coord}}(d_{\mathrm{CE}}u_f)
  =X_f=d_{\pi_{\mathrm{lin}}}O_f.
\]
Finite Taylor-degree maps in this construction are coefficient
projections of the row-finite coordinate product.  All Lie brackets,
Chevalley--Eilenberg differentials, and Schouten operations are formed
in the completed algebra and then evaluated by these projections.

\begin{defn}[Bracket-admissible \(B_q\)]
\label{def:completion-bracket-admissible}
Fix \(q>1\).  Let
\(B_q\subset
\mathrm{PV}_{\mathrm{coord}}(A_{\partial,H}^{\mathrm{st}})\)
be the closure of the polynomial subalgebra
\[
  \C[O_f,\partial_\rho,X_f\,:\,
    f\in\mathfrak h_q,\ \rho\in P_q],
  \qquad X_f=d_{\pi_{\mathrm{lin}}}O_f,
\]
in the weighted K\"othe topology of
Definition~\ref{def:completion-weighted-kothe}.  Then \(B_q\) is
complete Hausdorff, closed under multiplication, closed under the
Schouten bracket
\([B_q,B_q]_{\mathrm{Sch}}\subset B_q\), and closed under the
Lichnerowicz differential
\(d_{\pi_{\mathrm{lin}}}B_q\subset B_q\).  Its generator relations
include
\[
  [\partial_\rho,O_f]_{\mathrm{Sch}}
  =\langle\rho,f\rangle,\qquad
  [O_f,O_g]_{\mathrm{Sch}}=0,\qquad
  d_{\pi_{\mathrm{lin}}}O_f=X_f.
\]
The degree-zero Hamiltonian boundary bracket is the separate relation
\(\{O_f,O_g\}_\partial=O_{\{f,g\}}\).  The continuous inclusion
\(B_q\hookrightarrow\mathrm{PV}_{\mathrm{coord}}\) detects the
coordinate generators.
\end{defn}

\begin{thm}[\(P_0\)-bracket lift on \(B_q\)]
\label{thm:completion-p0-lift}
The regraded CE/PV cochain map
\(\Phi_{\mathrm{coord}}\) of
Theorem~\ref{thm:formal-disk-completed-ce-pv} restricts to a
\(P_0\)-isomorphism
\(\Phi^{B_q}_{P_0}:\Phi^{-1}_{\mathrm{coord}}(B_q)\xrightarrow{\sim}B_q\)
onto the bracket-admissible \(B_q\).  Its standard CE source has
\(|c^\rho|_{\mathrm{CE}}=1\) and \(|u_f|_{\mathrm{CE}}=2\);
\(\operatorname{Regr}_u\) gives the displayed coordinate degree
\(|u_f|_{\mathrm{coord}}=0\).
\end{thm}

\begin{proof}
By Definition~\ref{def:completion-bracket-admissible},
\(B_q\) is complete Hausdorff, multiplicatively closed,
Schouten-closed, \(d_{\pi_{\mathrm{lin}}}\)-closed, and contains the
coordinate generators \(O_f,\partial_\rho,X_f\).
Each of the five conditions in
Definition~\ref{defn:bracket-kernel-admissible} is satisfied; the
\(P_0\)-bracket lift criterion of
Theorem~\ref{thm:universal-ce-pv-koszul-criterion}(2) applies and
produces the isomorphism.
\end{proof}

\begin{defn}[Kernel-admissible completion \(\mathcal K_{B_q}\)]
\label{def:completion-kernel-admissible}
Let \(\mathcal K_{B_q}=\mathrm{Conv}^{\mathrm{cts}}_{\mathsf C}
\bigl(C^{\mathrm{CE}}_\bullet(\mathfrak g),B_q\bigr)\) be the
continuous convolution dg Lie algebra of CE chains of
\(\mathfrak g=\mathfrak h_q\ltimes\mathcal P^\Pi_q[1]\) into
\(B_q\), with bracket and differential induced by the operad
structure on \(B_q\).  The diagonal tensor
is defined from coefficient-dual families
\(\{H_I\}\subset\mathfrak h_q\) and
\(\{\rho^I\}\subset P_q\).  Write
\(O_I=O_{H_I}\), \(\partial^I=\partial_{\rho^I}\), and
\(\eta_I=\eta_{\rho^I}\in P_q[1]\).  Then
\[
  \Theta_{B_q}
  =\sum_I H_I\otimes\partial^I
   +\sum_I\eta_I\otimes O_I
  \in\mathfrak g_q\widehat\otimes B_q
\]
by the K\"othe continuity
  estimate
  (Proposition~\ref{prop:completion-promatlis-continuity}), and the
  sum converges in \(\mathcal K_{B_q}\).  The corresponding
  \(\mathcal B^{\mathrm{adm}}\)-datum is obtained by adjoining
  \(\mathcal K_{B_q}\) to the kernel list and the convolution differential
  and bracket to the operation list.
\end{defn}

\begin{ex}[Two completion tests for the diagonal Casimir]
\label{ex:completion-toy-casimir}
Let
\[
  H_d=\operatorname{Span}\{z_1^az_2^b:a+b=d\},\qquad
  \mathfrak h=\prod_{d\ge1}H_d .
\]
This example is the coefficient test behind
Proposition~\ref{prop:verified-badm-weighted-model}.  A finite window
tests only the degree pieces \(H_d\) with \(d\le M\).  A completed
kernel exists exactly when these finite-window tensors form a compatible
inverse-limit family in the declared completed tensor product.  The
maps \(\bigoplus_dH_d\to H_{\leq M}\) are coefficient projections;
brackets are formed in the full Hamiltonian algebra and then
projected to these spaces.

For a finite window \(M\), put
\[
  H_{\le M}=\bigoplus_{1\le d\le M}H_d,\qquad
  H_{\le M}^{\vee}=\bigoplus_{1\le d\le M}H_d^\vee .
\]
The finite Casimir
\[
  C_{\le M}=\sum_{1\le d\le M}\sum_i e_{d,i}\otimes e^i_d
  \in H_{\le M}\otimes H_{\le M}^{\vee}
\]
is an honest tensor.  If \(M'\ge M\), the projection
\[
  H_{\le M'}\otimes H_{\le M'}^\vee
  \longrightarrow
  H_{\le M}\otimes H_{\le M}^\vee
\]
sends \(C_{\le M'}\) to \(C_{\le M}\).

For a degree weight \(w\) satisfying
Definition~\ref{defn:wt-degree-weight}, the weighted finite tensors
\[
  C_{\le M,w}
  =
  \sum_{1\le d\le M}\sum_i
  \bigl(w(d)e_{d,i}\bigr)\otimes
  \bigl(w(d)^{-1}e^i_d\bigr)
\]
are compatible under the truncation maps
\[
  H_{\le M}^{(w)}=\bigoplus_{d\le M}w(d)H_d,\qquad
  (H_{\le M}^{\vee})^{(w^{-1})}
    =\bigoplus_{d\le M}w(d)^{-1}H_d^\vee .
\]
For \(M'\ge M\), the truncation maps
\(H_{\le M'}^{(w)}\twoheadrightarrow H_{\le M}^{(w)}\) and
\((H_{\le M'}^\vee)^{(w^{-1})}\twoheadrightarrow
(H_{\le M}^\vee)^{(w^{-1})}\) satisfy
\[
  (\pi_{M'M}\otimes\pi^\vee_{M'M})(C_{\le M',w})=C_{\le M,w}.
\]
Their inverse-limit class is the convergent weighted Casimir
\[
  C_{\mathfrak h,w}\in
  \mathfrak h_w\widehat\otimes
  D^{\mathrm{BV}}_w
\]
of Proposition~\ref{prop:wt-casimir-convergence}.  This inverse system
is the test for the admissible completion: every projection is finite
dimensional, the inverse system is Mittag--Leffler, and the displayed
kernel is a continuous coevaluation in the declared completed tensor
product.  The word ``convergent'' here means convergence in this
inverse-limit topology.  It is not a claim that the same infinite tensor
converges in the strict product/direct-sum Tate pair.  The three tests
are therefore
\[
\begin{array}{rcl}
  H_{\le M}\otimes H_{\le M}^{\vee}
  &:&
  C_{\le M}\ \text{is an honest finite tensor},\\[2mm]
  \mathfrak h_w\widehat\otimes
  D^{\mathrm{BV}}_w
  &:&
  (C_{\le M,w})_M\ \text{has an inverse-limit value},\\[2mm]
  \mathfrak h_\Pi\widehat\otimes D_\oplus
  &:&
  (C_{\le M})_M\ \text{has no limit in this tensor product}.
\end{array}
\]

In the strict product/direct-sum Tate pair
\[
  \mathfrak h_\Pi=\prod_{d\ge1}H_d,\qquad
  D_\oplus=\bigoplus_{d\ge1}H_d^\vee ,
\]
the same formal expression
\[
  C_{\mathfrak h}=\sum_{d\ge1}\sum_i e_{d,i}\otimes e^i_d
\]
does not define an element of
\(\mathfrak h_\Pi\widehat\otimes D_\oplus\).  The second tensor factor
of any element has finite support in the \(D_\oplus\) degree direction,
whereas \(C_{\mathfrak h}\) has nonzero component in every
\(H_d^\vee\).  Equivalently, such a tensor has nonzero image under
the projection to
\(\mathfrak h_\Pi\widehat\otimes H_d^\vee\) for every \(d\), which
contradicts the finite degree support imposed by
\(D_\oplus=\bigoplus_d H_d^\vee\).  This is
Lemma~\ref{lem:tate-casimir-obstruction} in the smallest example.  The
strict pair has evaluation
\(\mathfrak h_\Pi\times D_\oplus\to\C\); it does not have the
coevaluation kernel required by the BV propagator.

Equivalently, let \(B_{\mathrm{toy},w}\) be the weighted closure of the
finite polynomial polyvectors generated by
\(\{O_{e_{d,i}},\theta_{e_{d,i}}\}_{d,i}\), closed under \(d_\pi\),
Schouten bracket, and pointwise product, and let
\(\mathcal K_{B_{\mathrm{toy},w}}\) be the diagonal tensor target
obtained from
\(\mathfrak h_w\widehat\otimes
D^{\mathrm{BV}}_w\).  The pair
\((B_{\mathrm{toy},w},\mathcal K_{B_{\mathrm{toy},w}})\) is the toy
bracket- and kernel-admissible datum: every finite coefficient of each
operation is obtained by projection to \(H_{\le M}\), compatibility
with the projection maps passes these coefficients to the
Mittag--Leffler inverse limit, and the Casimir is the element
\(C_{\mathfrak h,w}\).  Let \(B_{\mathrm{toy},\Pi}\) be the same
finite-window algebra placed in the strict product/direct-sum coefficient
pair \((\mathfrak h_\Pi,D_\oplus)\).  Its projected bracket
coefficients exist, while kernel-admissibility would require an element of
\(\mathfrak h_\Pi\widehat\otimes D_\oplus\) with the finite-window images
\((C_{\le M})_M\).  The absence of such an element makes
kernel-admissibility an additional condition.  This example supplies the coefficient
coevaluation test only; a Costello QME counterterm tower is the separate
analytic datum of Problem~\ref{prob:weighted-rg-locality}.
\end{ex}

\begin{thm}[Maurer--Cartan kernel]
\label{thm:completion-mc-kernel}
The diagonal tensor \(\Theta_{B_q}\) satisfies the Maurer--Cartan
equation
\(d\Theta_{B_q}+\tfrac12[\Theta_{B_q},\Theta_{B_q}]_{\mathcal K_{B_q}}=0\)
in the kernel-admissible completion \(\mathcal K_{B_q}\) of
Definition~\ref{def:completion-kernel-admissible}.
\end{thm}

\begin{proof}
The coordinate identities
\[
  [\partial_\rho,O_f]_{\mathrm{Sch}}
    =\langle\rho,f\rangle,
  \qquad
  d_{\pi_{\mathrm{lin}}}O_f=X_f,
  \qquad
  d_{\pi_{\mathrm{lin}}}\partial_\rho
    =\Phi_{\mathrm{coord}}(d_{\mathrm{CE}}c^\rho)
\]
identify the Maurer--Cartan equation for
\(\Theta_{B_q}\) with the Chevalley--Eilenberg differential of
\(\mathfrak g_q\).  Continuity of the bracket
(Proposition~\ref{prop:completion-promatlis-continuity}) makes the
infinite sum converge, and the Maurer--Cartan equation reduces to
the Jacobi identity in \(\mathfrak h_q\), the coadjoint-module
identity on \(P_q\), and the operadic \(P_0\)-compatibility on
\(B_q\).
\end{proof}

\subsection{The Tate bar-cobar admissible envelope}
\label{ssec:completion-tate-bar-cobar}

Level 4 of the CE/PV ladder is the bar-cobar comparison
\(\Omega C^{\mathrm{CE}}_\bullet(\mathfrak g_{\mathrm{adm}})\xrightarrow{\sim}
U(\mathfrak g_{\mathrm{adm}})\) for an admissible Lie algebra
replacement.

\begin{defn}[Admissible Tate envelope]
\label{def:completion-admissible-tate}
An admissible Tate envelope datum for
\(\mathfrak g=\mathfrak h_q\ltimes\mathcal P^\Pi_q[1]\) consists of
a strict inverse system of finite-dimensional nilpotent dg Lie
quotients
\[
  q_M\colon\mathfrak g\longrightarrow\mathfrak g_M,
  \qquad
  \mathfrak g_{M+1}\twoheadrightarrow\mathfrak g_M,
\]
a continuous comparison
\(\mathfrak g\to\mathfrak g_{\mathrm{adm}}
:=\varprojlim_M\mathfrak g_M\), and compatible conilpotent,
Poincar\'e--Birkhoff--Witt, and Roos filtrations on
\[
  C^{\mathrm{CE}}_\bullet(\mathfrak g_{\mathrm{adm}})
  :=\mathbf R\!\varprojlim_M
    C^{\mathrm{CE}}_\bullet(\mathfrak g_M),
  \qquad
  U(\mathfrak g_{\mathrm{adm}})
  :=\mathbf R\!\varprojlim_M U(\mathfrak g_M).
\]
The Taylor maps
\(\mathfrak h_q\to\mathfrak m/\mathfrak m^{W+1}\) enter this
datum solely as coefficient projections for coordinate formulas.
The dg Lie quotient maps \(q_M\) and their transition homotopies are
the supplied pro-Lie part of the envelope datum.
\end{defn}

\begin{lem}[Milnor exact sequence for \(C^{\mathrm{CE}}_\bullet\)]
\label{lem:completion-milnor-ce}
For the Tate envelope \(\mathfrak g_{\mathrm{adm}}\) of
Definition~\ref{def:completion-admissible-tate},
\[
  0\to\varprojlim_M{}^1H^{\bullet-1}(C^{\mathrm{CE}}_\bullet(\mathfrak g_M))
   \to H^\bullet(C^{\mathrm{CE}}_\bullet(\mathfrak g_{\mathrm{adm}}))
   \to\varprojlim_M H^\bullet(C^{\mathrm{CE}}_\bullet(\mathfrak g_M))
   \to 0,
\]
for every cohomological degree in which the sequential Roos
totalization is defined.  A degreewise Mittag--Leffler cohomology
system has vanishing \(\varprojlim^1\)-term in that degree.
\end{lem}

\begin{proof}
The two-column Roos resolution for the sequential inverse system of
Chevalley--Eilenberg chain complexes gives the displayed Milnor exact
sequence by axiom C1.  Degreewise Mittag--Leffler on cohomology
annihilates the corresponding derived-limit term.
\end{proof}

\begin{thm}[Tate bar-cobar comparison with derived-limit obstruction]
\label{thm:completion-tate-bar-cobar}
At each finite quotient the map
\[
\Omega C^{\mathrm{CE}}_\bullet(\mathfrak g_M)\to U(\mathfrak g_M)
\]
is a quasi-isomorphism.  Suppose the admissible Tate envelope datum
also supplies a filtered comparison
\[
  \kappa_\Omega\colon
  \Omega\mathbf R\!\varprojlim_M
    C^{\mathrm{CE}}_\bullet(\mathfrak g_M)
  \longrightarrow
  \mathbf R\!\varprojlim_M
    \Omega C^{\mathrm{CE}}_\bullet(\mathfrak g_M)
\]
which is a quasi-isomorphism, together with degreewise
Mittag--Leffler vanishing for the finite bar--cobar cones.  Then the
completed map
\(\Omega C^{\mathrm{CE}}_\bullet(\mathfrak g_{\mathrm{adm}})\to
U(\mathfrak g_{\mathrm{adm}})\) is a quasi-isomorphism in
\(\mathcal B^{\mathrm{adm}}\), through the presentation \(\mathsf C\).
Its general obstruction is the cohomology of the Roos-totalized
bar--cobar cone together with the cone of \(\kappa_\Omega\).
\end{thm}

\begin{proof}
Each \(\mathfrak g_M\) is finite nilpotent, hence the bar-cobar
comparison \(\Omega C^{\mathrm{CE}}_\bullet(\mathfrak g_M)\to
U(\mathfrak g_M)\) is a quasi-isomorphism of dg algebras
\cite{quillen-htpy}\cite{kontsevich-dq}.
Apply the Roos functor to the finite bar--cobar cones.  Their
degreewise Mittag--Leffler hypothesis makes the totalized cone
acyclic.  Composing this totalized quasi-isomorphism with
\(\kappa_\Omega\) gives the completed comparison in the filtered
topology presented by \(\mathsf C\), hence in the corresponding
\(\mathcal B^{\mathrm{adm}}\)-object.  The two displayed cones record
the residual obstruction for any other pro-Lie system.
\end{proof}

\subsection{\texorpdfstring{Open-trace \(A_\infty\)-Koszul homotopy datum}{Open-trace A--Koszul homotopy datum}}
\label{ssec:completion-ainfty-koszul}

Level 5 of the CE/PV ladder requires an
\((A_{\mathrm{op}},C_{\mathrm{op}},B_{\mathrm{CE/PV}},\tau_{\mathrm{op}})\)
datum.

\begin{defn}[Open-trace \(A_\infty\)-Koszul datum]
\label{def:completion-open-ainfty-koszul}
The open-trace filtered cyclic \(A_\infty\) algebra is the
non-effective fixed-rank representative
\[
  A_{\mathrm{op}}=A^{\mathrm{cl}}_{\partial,N}
  =
  C^\bullet_{\mathrm{CE}}(\mathfrak{gl}_N,\mathrm{Kosz}_N)
\]
of Definition~\ref{def:E1-brane-algebra}.  The effective Hamiltonian
representative replaces the odd \(\mathfrak{gl}_N[1]\)-coefficient in
\(\mathrm{Kosz}_N\) by \(\mathfrak{sl}_N[1]\) and the gauge Lie algebra by
\(\mathfrak{pgl}_N\); the trace-sector reduction applies after the central
\(\operatorname{Tr}\psi\) excess is removed.  The cyclic structure is
the supplied map
\(\Tr_{\mathrm{open}}:A_{\mathrm{op}}\to\C[-1]\).  Its ghost-zero
closed-to-open trace observables at fixed rank are
\[
  J_N^s(f)=
  \bigl[\Tr s(f)(\phi_1,\phi_2)\bigr],
  \qquad f\in\bar A,
\]
with the finite-\(N\) trace relations retained.  Effective scalar
reduction removes \(\operatorname{Tr}\psi\) and the empty trace before
adjoining the unit of the reduced observable algebra.  Degreewise
stable passage gives
\(J^{\mathrm{st}}\colon\mathfrak h_q\to
A_{\partial,H}^{\mathrm{st}}\).  For each Taylor degree \(W\), the
map \(J_{\leq W}^{\mathrm{st}}\) is the coefficient projection of
this stable trace map.  The Koszul-dual conilpotent filtered
coalgebra is
\(C_{\mathrm{op}}=BA_{\mathrm{op}}=
T^c(s\overline{A}_{\mathrm{op}})\) (the bar construction), with the
filtration by tensor length.  The bracket-admissible CE/PV target
is \(B_{\mathrm{CE/PV}}=B_q\) of
Definition~\ref{def:completion-bracket-admissible}.  The twisting
cochain is
\(\tau_{\mathrm{op}}:C_{\mathrm{op}}\to A_{\mathrm{op}}\), defined
by the canonical length-one bar twisting cochain together with the
cyclic primitive comparison of
Loday--Quillen \cite{loday-quillen}*{Thm.~6.2} and Tsygan
\cite{tsygan}*{Thm.~1}.
\end{defn}

\begin{thm}[Twisting cochain \(\tau_{\mathrm{op}}\) satisfies MC]
\label{thm:completion-tau-op-mc}
The twisting cochain \(\tau_{\mathrm{op}}\) of
Definition~\ref{def:completion-open-ainfty-koszul} satisfies the
Maurer--Cartan equation
\[
  d\tau_{\mathrm{op}}
  +\mu_2(\tau_{\mathrm{op}}\otimes\tau_{\mathrm{op}})\Delta_2
  +\sum_{k\ge 3}\mu_k(\tau_{\mathrm{op}}^{\otimes k})\Delta_k=0,
\]
and the bar-cobar comparison
\(\Omega C_{\mathrm{op}}\xrightarrow{\sim}A_{\mathrm{op}}\) is a
filtered cyclic \(A_\infty\) quasi-isomorphism.  At fixed rank, its
primitive connected ghost-zero image is the span of \(J_N^s(f)\)
modulo the finite-\(N\) trace relations.  After effective scalar
reduction and degreewise stable passage, this image is
\[
  H^0(A_{\mathrm{op}}^{\mathrm{st}})_{\mathrm{prim,conn}}
  \xrightarrow{\sim}
  \mathrm{span}\{J^{\mathrm{st}}(f)\,:\,
    f\in\mathfrak h_q\}.
\]
\end{thm}

\begin{proof}
Write the \(A_\infty\) structure maps of \(A_{\mathrm{op}}\) as
\(\mu_k:A_{\mathrm{op}}^{\otimes k}\to A_{\mathrm{op}}[2-k]\).  The
bar differential on
\[
  C_{\mathrm{op}}=BA_{\mathrm{op}}=T^c(s\overline A_{\mathrm{op}})
\]
is the coderivation
\[
  b=\sum_{k\ge1} b_k,\qquad
  b_k=s\,\mu_k\,s^{-k}
\]
applied to consecutive \(k\)-blocks.  The cyclic \(A_\infty\) identities
are \(b^2=0\), together with invariance of the cyclic trace under each
\(\mu_k\).  The twisting cochain is the projection to the length-one
summand followed by desuspension,
\[
  \tau_{\mathrm{op}}\colon
  T^c(s\overline A_{\mathrm{op}})
  \twoheadrightarrow s\overline A_{\mathrm{op}}
  \xrightarrow{s^{-1}}\overline A_{\mathrm{op}}
  \hookrightarrow A_{\mathrm{op}} .
\]
Taking the length-one component of \(b^2=0\) gives exactly
\[
  d\tau_{\mathrm{op}}
  +\mu_2(\tau_{\mathrm{op}}\otimes\tau_{\mathrm{op}})\Delta_2
  +\sum_{k\ge3}\mu_k(\tau_{\mathrm{op}}^{\otimes k})\Delta_k=0,
\]
with the Koszul signs fixed by the suspension convention.  Hence
\(\tau_{\mathrm{op}}\) is a Maurer--Cartan element of the convolution dg
Lie algebra
\(\mathrm{Conv}(C_{\mathrm{op}},A_{\mathrm{op}})\).

The counit
\[
  \epsilon\colon \Omega C_{\mathrm{op}}\longrightarrow A_{\mathrm{op}}
\]
is filtered by tensor length.  On the associated graded, only the
linear differential \(\mu_1\) and the universal bar--cobar contraction
remain.  The cone of \(\operatorname{gr}\epsilon\) decomposes into the
positive tensor-length summands of the augmented tensor coalgebra; the
extra degeneracy inserting the coaugmentation gives a contracting
homotopy on every summand of length \(>1\).  Therefore
\(\operatorname{gr}\epsilon\) is a quasi-isomorphism.  The filtration is
exhaustive, complete, and bounded below in each finite tensor window, so
the spectral sequence comparison theorem gives that \(\epsilon\) is a
filtered \(A_\infty\) quasi-isomorphism.  Cyclicity is preserved because
the bar coderivation components \(b_k\) are the suspended cyclic
operations \(\mu_k\), and the contraction uses the coaugmentation, which
is orthogonal to the augmentation ideal for the declared cyclic pairing.

Finally, in each degreewise stable trace range the
Loday--Quillen--Tsygan primitive
projection identifies the primitive connected summand with single cyclic
words.  The Koszul differential \(Q\psi=[\phi_1,\phi_2]\) kills adjacent
transpositions by the homotopy
\(Q\operatorname{Tr}(\psi vu)=
\operatorname{Tr}(u\phi_1\phi_2v)-
\operatorname{Tr}(u\phi_2\phi_1v)\).  Thus every primitive connected
class has the unique commutative representative
\(\overline{\operatorname{Tr} f(\phi_1,\phi_2)}
=J^{\mathrm{st}}(f)\) in the stable trace system.  At fixed rank the
same argument lands in the finite trace quotient; compatible
degreewise passage gives the displayed stable identification.
\end{proof}

\subsection{\texorpdfstring{Hochschild equivariance homotopies for \(\Phi_N^{\mathrm{Hoch}}\)}{Hochschild equivariance homotopies for Phi N Hoch}}
\label{ssec:completion-hochschild-centrality}

The stable coordinate map has standard Chevalley--Eilenberg source
degrees \(|c^\rho|_{\mathrm{CE}}=1\) and
\(|u_f|_{\mathrm{CE}}=2\), followed by the coordinate regrading
\(\operatorname{Regr}_u\):
\[
  c^\rho\longmapsto\partial_\rho,\qquad
  u_f\longmapsto O_f=J^{\mathrm{st}}(f),\qquad
  d_{\mathrm{CE}}u_f\longmapsto
  X_f=d_{\pi_{\mathrm{lin}}}O_f.
\]
Its Hochschild realization factors through two distinct comparisons.
The ordinary Hochschild--Kostant--Rosenberg map realizes the
zero-Poisson special fibre.  The linear-Poisson differential is
realized by an admissible continuous filtered twisted-formality,
Poincar\'e--Birkhoff--Witt, and Roos datum.  A fixed-rank morphism
\(\Phi_N^{\mathrm{Hoch}}\colon
\mathcal A^{\mathrm{cl}}_{\mathrm{bulk},0}\to
Z^{\mathrm{der}}_{E_1}(A^{\mathrm{cl}}_{\partial,N})\) is a
separate Maurer--Cartan lift extending the fixed-rank trace
cochains and Hamiltonian derivations.

\begin{prop}[Zero-Poisson Hochschild--Kostant--Rosenberg target]
\label{prop:completion-hkr-zero-poisson}
Let \(V\subset\mathfrak h_q\) be a finite-dimensional coordinate
subspace and \(A_V=\mathbf S(V)\).  The
Hochschild--Kostant--Rosenberg map is a quasi-isomorphism
\[
  \operatorname{HKR}_V\colon
  (\mathrm{PV}(A_V),0)
  \longrightarrow
  (C^\bullet_{\mathrm{Hoch}}(A_V,A_V),b).
\]
A cofinal system of finite coordinate subspaces, transition maps,
and Hochschild--Kostant--Rosenberg-compatible transition homotopies
has a Roos totalization
\[
  \operatorname{HKR}^{\mathrm{cont}}_{\mathrm{adm},0}\colon
  (\mathrm{PV}_{\mathrm{adm}}
    (A_{\partial,H}^{\mathrm{st}}),0)
  \xrightarrow{\sim}
  C^\bullet_{\mathrm{Hoch},\mathrm{adm},0}
    (A_{\partial,H}^{\mathrm{st}},
     A_{\partial,H}^{\mathrm{st}}).
\]
The finite subspaces and Taylor-degree windows in this display are
coefficient projections.  Their transition homotopies are part of
the supplied Roos datum.
\end{prop}

\begin{proof}
The finite algebra \(A_V\) is smooth and polynomial, so the ordinary
Hochschild--Kostant--Rosenberg theorem gives the first
quasi-isomorphism.  Applying the Roos totalization to the compatible
system and its transition homotopies gives the completed special
fibre.
\end{proof}

\begin{defn}[Centre-map cochains and homotopies]
\label{def:completion-centre-cochains}
Put \(A=A_{\partial,H}^{\mathrm{st}}\).  Supply an admissible twisted
formality and Poincar\'e--Birkhoff--Witt datum
\[
  \mathcal D_{\mathrm{form}}
  = (\mathcal U,m_\hbar,\beta_{\mathrm{PBW}},
     \{H^{\mathrm{Roos}}_{\alpha\beta}\})
\]
for \((A,\pi_{\mathrm{lin}})\) in the sense of
Definition~\ref{defn:admissible-twisted-formality-pbw}.  Its twisted
map, composed with the regraded coordinate isomorphism, is
\[
  \Psi_\hbar
  =
  \mathcal U^{\hbar\pi}\circ\Phi_{\mathrm{coord}}\colon
  \bigl(\operatorname{Regr}_u
    C^\bullet_{\mathrm{CE},\mathrm{coord}}(\mathfrak g_q)
    [[\hbar_{\mathrm{Rees}}]],
    \hbar_{\mathrm{Rees}}d_{\mathrm{CE}}\bigr)
  \longrightarrow
  \bigl(C^\bullet_{\mathrm{Hoch},\mathrm{adm}}
    (A_\hbar,A_\hbar),b_\hbar\bigr),
\]
where
\[
  A_\hbar=(A[[\hbar_{\mathrm{Rees}}]],m_\hbar),
  \qquad
  \beta_{\mathrm{PBW}}\colon A_\hbar
  \xrightarrow{\sim}
  \widehat U_{\hbar_{\mathrm{Rees}}}(\mathfrak h_q).
\]
Its generator boundary values are
\[
  \Psi_\hbar(c^\rho)
    =\operatorname{HKR}(\partial_\rho)
      +O(\hbar_{\mathrm{Rees}}),
  \qquad
  \Psi_\hbar(u_f)
    =M_{O_f}+O(\hbar_{\mathrm{Rees}}),
\]
and
\[
  \Psi_\hbar(d_{\mathrm{CE}}u_f)
    =\operatorname{HKR}(X_f)
      +O(\hbar_{\mathrm{Rees}}),
  \qquad
  b_\hbar\Psi_\hbar(u_f)
    =\hbar_{\mathrm{Rees}}
      \Psi_\hbar(d_{\mathrm{CE}}u_f).
\]

At fixed rank, prescribe the boundary cochains
\[
  M_{J_N^s(f)}\in
  CC^0(A_{\mathrm{op}},A_{\mathrm{op}}),
  \qquad
  \Theta^{(1)}_{f,N}\in
  CC^1(A_{\mathrm{op}},A_{\mathrm{op}}),
\]
where \(\Theta^{(1)}_{f,N}\) represents the chosen Koszul
Hamiltonian derivation \(X_{f,N}\).  A finite-rank centre map is a
Maurer--Cartan solution extending these cochains, the finite-rank
trace relations, and the brane-current cochains
\(T_\rho\in CC^\bullet(A_{\mathrm{op}},A_{\mathrm{op}})\).
Its product, \(P_0\), and Jacobiator homotopies are denoted
\(H^{\mathrm{prod}}_{x,a},H^{P_0}_{x,a},H_3,\ldots\).
\end{defn}

\begin{thm}[Hochschild equivariance criterion on the bracket-admissible sector]
\label{thm:completion-centrality-homotopies}
A supplied datum \(\mathcal D_{\mathrm{form}}\) gives the stable
linear-Poisson Hochschild comparison \(\Psi_\hbar\) of
Definition~\ref{def:completion-centre-cochains}.  A fixed-rank
Hochschild-centre lift extending
\(\Theta^{(1)}_{f,N},M_{J_N^s(f)}\) exists precisely when there is a
continuous assignment of these Hochschild cochains and homotopies
\(H_{f,g},H_{f|g},H^{\mathrm{mult}}_{f,g}\), compatible with the
weighted K\"othe filtration, satisfying
\[
  \begin{aligned}
    [\Theta^{(1)}_{f,N},\Theta^{(1)}_{g,N}]_G
      &=\Theta^{(1)}_{\{f,g\},N}+b_NH_{f,g},\\
    [\Theta^{(1)}_{f,N},M_{J_N^s(g)}]_G
      &=M_{J_N^s(\{f,g\})}+b_NH_{f|g},\\
    [M_{J_N^s(f)},M_{J_N^s(g)}]_G
      &=b_NH^{\mathrm{mult}}_{f,g},
  \end{aligned}
\]
for all \(f,g\in\mathfrak h_q\), together with the corresponding
Jacobiator homotopy \(H_3\) and the finite-rank trace relations.
For such a lift the obstruction component
\(\mathrm{ob}^{\mathrm{prod}}_{\mathrm{cent}}\in
H^2(\mathfrak K_{\mathrm{cent}})\) of the four-component
\(\operatorname{Ob}_{\mathrm{cent/QME/desc}}\) vanishes on the
bracket-admissible sector.  Conversely, vanishing of this component is
the obstruction-theoretic statement that such Hochschild homotopies may
be chosen in the selected Hochschild model.
The formality extension classes
\(\operatorname{Ob}_{\mathcal U,m+1}\), the PBW comparison cone,
and the Roos transition defects govern the construction of
\(\mathcal D_{\mathrm{form}}\).  The \(P_0\)-centrality,
brane-defect quantum master equation, and centre-map Roos descent
entries govern the subsequent finite-rank lift.
\end{thm}

\begin{proof}
Proposition~\ref{prop:completion-hkr-zero-poisson} gives the ordinary
special fibre.  Definition~\ref{defn:admissible-twisted-formality-pbw}
twists the supplied filtered \(L_\infty\) quasi-isomorphism by
\(\hbar_{\mathrm{Rees}}\pi_{\mathrm{lin}}\), transports the resulting
star product through \(\beta_{\mathrm{PBW}}\), and totalizes the
transition homotopies.  This gives the stable comparison
\(\Psi_\hbar\).

The definition of \(\mathfrak K_{\mathrm{cent}}\) is the mapping cone
for the Gerstenhaber product, the \(P_0\)-bracket, and the first
Jacobiator of the candidate centre map.  The displayed equations are
the degree-two components of its Maurer--Cartan equation after the
stable PV identities have been transported through
\(\Psi_\hbar\) and evaluated against the prescribed fixed-rank
boundary cochains.  Thus a choice of
\(H_{f,g},H_{f|g},H^{\mathrm{mult}}_{f,g}\)
and \(H_3\) is precisely a primitive for
\(\mathrm{ob}^{\mathrm{prod}}_{\mathrm{cent}}\).  The formality,
PBW, and Roos cones occur one stage earlier and supply the Hochschild
model in which this primitive is sought.
\end{proof}

\begin{rmk}[Cotangent centrality and \(P_0\)-bracket lifts]
\label{rmk:completion-cotangent-centrality}
The constant Schouten field \(\partial_\rho\) is the coordinate image
of \(c^\rho\).  Its fixed-rank current realization is supplied by a
principal-part cochain
\(T_\rho\in
CC^\bullet(\mathrm{Obs}_{\partial,N},
           \mathrm{Obs}_{\partial,N})\)
and the brane-defect propagator extension of
\S\ref{ssec:completion-graph-qme}.  The corresponding product and
\(P_0\)-centrality homotopies are entries of
\(\operatorname{Ob}_{\mathrm{cent/QME/desc}}\).
\end{rmk}

\subsection{Protected-sector cyclic BV strong deformation retract}
\label{ssec:completion-protected-sector-sdr}

The local theorem on \(\R^2\times\C^2\) is presented as a formal
Hamiltonian stalk; for the universal completed theorem the local
fields must arise from a parent BV/factorization target by an actual
deformation retract.

\begin{defn}[Admissible protected BV/factorization target]
\label{def:completion-admissible-protected-target}
An admissible protected mixed holomorphic-topological BV/factorization
target is a tuple
\(\mathcal T=(Y,\Omega_Y,\mathcal B,p,\mathcal E_T,S_T,\omega_T,
\mathcal Q,\mathcal C_{\mathcal T})\) where \(Y\) is the target
geometry, \(\Omega_Y\) is the holomorphic-symplectic volume datum,
\(\mathcal B\) is the brane category, \(p\in Y\) is the brane
point, \((\mathcal E_T,Q_T,\omega_T,S_T)\) is the parent BV
theory, \(\mathcal Q\) is the protected supercharge or BRST
operator, and \(\mathcal C_{\mathcal T}\) is the matched-conventions
datum (\S\ref{app:comparison-conventions}).
\end{defn}

\begin{defn}[Protected-sector cyclic BV strong deformation retract]
\label{def:completion-protected-sector-sdr}
A protected-sector cyclic BV strong deformation retract for a
target \(\mathcal T\) is a triple of \(Q_T\)-equivariant continuous
chain maps and homotopy
\((i_T,p_T,h_T)\):
\(i_T:\mathcal E_H\hookrightarrow\mathcal E_T\),
\(p_T:\mathcal E_T\twoheadrightarrow\mathcal E_H\),
\(h_T:\mathcal E_T\to\mathcal E_T[-1]\), satisfying
\[
  p_Ti_T=\mathrm{id}_{\mathcal E_H},
  \qquad
  Q_Th_T+h_TQ_T=\mathrm{id}_{\mathcal E_T}-i_Tp_T,
  \qquad
  i_T^*\omega_T=\omega_H,
  \qquad
  h_T^\vee=\pm h_T
\]
relative to the BV pairing.  The transferred interaction is
\(S_T^{\mathrm{trans}}=S_H+\sum_{k\ge 3}S^{\mathrm{trans}}_k\), with
the \(L_\infty\)-transferred higher operations
\(S^{\mathrm{trans}}_k\) computed by the Homological Perturbation
Lemma applied to \((i_T,p_T,h_T)\).
\end{defn}

\begin{thm}[Protected-sector obstruction]
\label{thm:completion-prot-sector-obstruction}
For an admissible protected target \(\mathcal T\) and a
protected-sector strong deformation retract (SDR) \((i_T,p_T,h_T)\),
the protected-sector obstruction
\[
  \mathrm{Ob}_{\mathrm{prot}}(\mathcal T)
  =\Bigl[\sum_{k\ge 3}S^{\mathrm{trans}}_k\Bigr]
  \in H^1\bigl(\mathfrak K_{\mathrm{prot}}(\mathcal T)\bigr),
  \qquad
  \mathfrak K_{\mathrm{prot}}(\mathcal T)
  =\mathrm{Def}_{\mathrm{cycBV}}(\mathcal E_H\rightleftarrows\mathcal E_T),
\]
is the curvature of the transferred interaction in the cyclic BV
deformation complex.  Strict universality
\(\mathrm{Ob}_{\mathrm{prot}}(\mathcal T)=0\) holds when the higher
\(S^{\mathrm{trans}}_k\) vanish; counterterm universality
\(\mathrm{Ob}_{\mathrm{prot}}(\mathcal T)=d_{\mathfrak K}\tau_{\mathrm{prot}}\)
holds when they are exact in the local-functional complex; projective
universality holds when they admit a central representative.
\end{thm}

\begin{proof}
The Homological Perturbation Lemma
\cite{crainic-perturbation-lemma}\cite{brown-perturbation-lemma} applied
to the SDR \((i_T,p_T,h_T)\) produces \(S^{\mathrm{trans}}_k\) by
the explicit formula
\(S^{\mathrm{trans}}_k=p_T\circ S_T\circ\cdots\circ S_T\circ h_T\circ S_T\circ i_T\)
with \(k-1\) propagators \(h_T\) interleaved by \(S_T\).  The
cyclicity hypothesis \(h_T^\vee=\pm h_T\) implies
\(S^{\mathrm{trans}}_k\) is cyclic.  The deformation complex
\(\mathfrak K_{\mathrm{prot}}(\mathcal T)\) records the obstruction for
the transferred interaction to equal \(S_H\); its first cohomology
class is the obstruction.  The three universality cases follow
from the obstruction-twist dichotomy
(Theorem~\ref{thm:obstruction-twist-dichotomy}) applied to
\(\mathfrak K_{\mathrm{prot}}\).
\end{proof}

\begin{rmk}[Existence of the protected retract]
\label{rmk:completion-protected-existence}
Theorem~\ref{thm:completion-prot-sector-obstruction} gives the
obstruction calculus once an SDR exists.  For a specific compact
Calabi--Yau threefold target an admissible protected SDR is part of the
compact field theory and is constrained by the matched-conventions
obstruction (\S\ref{app:comparison-conventions}).  The local theorem is the
formal Hamiltonian normal-slice test of the bulk produced by such an
SDR.
\end{rmk}

\subsection{Scalar-contact projection and the QME counterterm tower}
\label{ssec:completion-qme-tower}

The reduced non-scalar Costello QME class
\([\operatorname{Ob}^{\mathrm{red}}_{w,\partial,\hbar_{\mathrm{QME}}}]\) is defined
on \(\ker\sigma_{\mathrm{sc}}\) for a filtered chain map
\(\sigma_{\mathrm{sc}}\), refining the associated graded scalar
symbol.  The statements below identify the obstruction tower for this
lift: \(\sigma_{\mathrm{sc}}\) exists after the scalar-lift classes
vanish, and the counterterm tower exists as a formal inverse-limit
family after the corresponding Milnor pairs vanish.  The \(\theta_3\)
row is a finite-window primitive calculation inside this criterion.

\begin{defn}[Scalar-contact filtration and projection]
\label{def:completion-scalar-contact}
The scalar-contact filtration
\(F^\bullet\mathfrak Q^\bullet_{w,\partial,\hbar_{\mathrm{QME}}}(I)\) is generated
by the powers of the scalar central span \(\C\cdot\mathbf 1\subset
\mathfrak h_{\hbar_{\mathrm{QME}},N}\): \(F^r\) consists of local functionals
whose expansion in the Capelli generators contains at least \(r\)
copies of the scalar.  The associated graded is
\(\mathrm{gr}^\bullet_F\mathfrak Q^\bullet_{w,\partial,\hbar_{\mathrm{QME}}}(I)\),
with associated-graded scalar symbol
\(\sigma_{\mathrm{sc}}^{\mathrm{gr}}:\mathrm{gr}^\bullet_F
\mathfrak Q^\bullet_{w,\partial,\hbar_{\mathrm{QME}}}(I)\to
C^\bullet_{\mathrm{Lie}}(\bar A;\C\gamma_1)[1][[\hbar_{\mathrm{QME}}]]\).  The
filtered chain map is
\[
  \sigma_{\mathrm{sc}}:\mathfrak Q^\bullet_{w,\partial,\hbar_{\mathrm{QME}}}(I)
   \longrightarrow C^\bullet_{\mathrm{Lie}}(\bar A;\C\gamma_1)[1][[\hbar_{\mathrm{QME}}]],
\]
with \(\mathrm{gr}_F(\sigma_{\mathrm{sc}})=\sigma_{\mathrm{sc}}^{\mathrm{gr}}\).
\end{defn}

\begin{thm}[Filtered chain map \(\sigma_{\mathrm{sc}}\) and obstruction tower]
\label{thm:completion-sigma-sc}
On an unreduced polynomial or balanced branch with a fixed chain scalar
projection, work in a separated complete filtration for which the
order-\(r\) correction lies in \(F^r\).  The associated-graded scalar
symbol \(\sigma_{\mathrm{sc}}^{\mathrm{gr}}\) lifts to a filtered chain
map \(\sigma_{\mathrm{sc}}\) inside this fixed filtered complex if and
only if the scalar-lift
obstruction tower
\(o^{(r)}_{\sigma,w}(I)\in H^1\bigl(
\Hom(\mathrm{gr}^r_F\mathfrak Q^\bullet_{w,\partial,\hbar_{\mathrm{QME}}}(I),
C^\bullet_{\mathrm{Lie}}(\bar A;\C\gamma_1)[1][[\hbar_{\mathrm{QME}}]])_{-r}\bigr)\)
vanishes for every \(r\ge 0\).  At order \(r=1\), the obstruction
\(o^{(1)}_{\sigma,w}\) is the Capelli scalar cocycle
\(\omega(f,g)=[\{f,g\}]_0\).  On the unreduced polynomial algebra
\(\mathfrak h_{\mathrm{poly}}=\C[z_1,z_2]\) it has primitive
\(\eta(f)=-[f]_0\) with \(d_{\mathrm{CE}}\eta=\omega\).  On the
reduced Hamiltonian algebra \(\bar A\) the class is
\([\bar c]\in H^2_{\mathrm{Lie}}(\bar A,\C)\) and is retained as a
projective central class.  At higher orders, the obstruction lies in
the iterated extension of the Capelli scalar span and is the class
defined by the preceding choices in the exact couple of the filtration.
In the
finite-window prototype \(\theta_3\) row, the displayed obstruction
is killed after adjoining a scalar-zero windowwise primitive
\((d_M C_{\theta_3,M}=-\mathsf E_{\theta_3,M},c=1)\), confirming
the finite-window enlargement in that row; the all-order case is the
separate assertion that the windowwise data extend
Milnor-compatibly across the inverse system of windows
(Theorem~\ref{thm:completion-milnor-obstruction-vector}).
\end{thm}

\begin{proof}
By Definition~\ref{def:completion-scalar-contact}, the obstruction
to lifting \(\sigma_{\mathrm{sc}}^{\mathrm{gr}}\) by one
filtration order is a class in
\(H^1\bigl(\Hom(\mathrm{gr}^r_F,C^\bullet_{\mathrm{Lie}})_{-r}\bigr)\).
At order \(r=1\) the class is \([\omega]\).  The primitive
\(\eta\) exists only on the unreduced algebra
(Lemma~\ref{lem:omega-cocycle}); after reduction the surviving class
is \([\bar c]\) and the target must be the central extension.  At
higher orders the class lives on the \(r\)-fold symmetric product of
the Capelli scalar span.  The sequence of corrections converges in the
filtered topology because the \(r\)-th correction lies in \(F^r\) and
the filtration is separated and complete
(Definition~\ref{def:completion-host-category}); without this
pronilpotence the statement is only finite order.  The all-order
existence reduces to the Milnor compatibility of the windowwise
counterterms across the inverse system of finite Hamiltonian
windows, controlled by
Theorem~\ref{thm:completion-milnor-obstruction-vector}.
\end{proof}

\begin{defn}[Residual non-scalar complex and obstruction]
\label{def:completion-residual-nonscalar}
The residual non-scalar complex is
\(\mathfrak Q^\bullet_{\mathrm{red}}(I)=\ker\sigma_{\mathrm{sc}}\)
of Theorem~\ref{thm:completion-sigma-sc}.  The reduced non-scalar
Costello QME class is
\([\operatorname{Ob}^{\mathrm{red}}_{w,\partial,\hbar_{\mathrm{QME}}}]\in
H^1(\mathfrak Q^\bullet_{\mathrm{red}}(I),d_{\mathrm{QME}})\),
defined as the curvature of the reduced interaction in the kernel
of the scalar-contact projection.
\end{defn}

\subsubsection{Finite-window QME counterterm recursion}
\label{sssec:completion-qme-tower-recursion}

\begin{defn}[Finite-window residual complex]
\label{def:completion-finite-window-residual}
For each Hamiltonian window \(M\) and each loop/order \(n\), let
\(Q^\bullet_{w,M}(I)=\ker(\sigma_{\mathrm{sc},M})\subset
\mathcal O^{\mathrm{bd}}_{\mathrm{loc}}(\mathcal E^M_w|_I;
A^{\mathrm{pp},w,M}_{\partial,\hbar_{\mathrm{QME}}}(I))\).  At order \(n\) the
window obstruction is \(\mathfrak o^{\mathrm{ns}}_{n,M}\in
Q^1_{w,M}(I)\).  A compatible counterterm system is
\(\{C_{n,M}\in Q^0_{w,M}(I)\}_{n,M}\) with
\(d_M C_{n,M}=-\mathfrak o^{\mathrm{ns}}_{n,M}\) and
\(\pi_{M,N}C_{n,M}=C_{n,N}\) for \(M\ge N\).
\end{defn}

\begin{thm}[Milnor obstruction vector and counterterm criterion]
\label{thm:completion-milnor-obstruction-vector}
Assume the finite-window residual complexes form an inverse system
closed under the boundary and bracket terms of the QME, the graph
weights are summable in the chosen weighted topology, and the residuals
\(\mathfrak o^{\mathrm{ns}}_{n,M}\) are compatible under truncation.
The obstruction to a compatible counterterm system at fixed order \(n\)
is the Milnor pair
\[
  \mathfrak D^{\mathrm{ns}}_n
  =\bigl(\,([\mathfrak o^{\mathrm{ns}}_{n,M}])_M,\,\lambda_n\,\bigr)
  \in\varprojlim_M H^1(Q^\bullet_{w,M}(I))
   \oplus\varprojlim_M{}^1H^0(Q^\bullet_{w,M}(I)),
\]
where the first component is the windowwise primitive class and
the second is the Milnor compatibility class.  Under the displayed
finite-to-limit hypotheses, an all-order QME counterterm tower as a
formal \(\hbar_{\mathrm{QME}}\)-adic inverse-limit family exists if and only if
\(\mathfrak D^{\mathrm{ns}}_n=0\) for every \(n\ge 1\); when this
holds the formal counterterm \(C^{\mathrm{red}}_{\hbar_{\mathrm{QME}},w}=
\sum_{n\ge 1}\hbar_{\mathrm{QME}}^n C_{n}\), with
\(C_n\in\mathbf R\!\varprojlim_M Q^0_{w,M}(I)\), is a formal
counterterm primitive for the QME curvature equation
\(d_{\mathrm{QME}}C^{\mathrm{red}}_{\hbar_{\mathrm{QME}},w}+
\tfrac12\{C^{\mathrm{red}}_{\hbar_{\mathrm{QME}},w},
C^{\mathrm{red}}_{\hbar_{\mathrm{QME}},w}\}+\cdots
=-\mathrm{Ob}^{\mathrm{red}}_{w,\partial,\hbar_{\mathrm{QME}}}\) at all orders as a
formal power series.  No analytic convergence in \(\hbar_{\mathrm{QME}}\) is asserted.
\end{thm}

\begin{proof}
The Milnor exact sequence (axiom C1) for the inverse system
\(\{Q^\bullet_{w,M}\}\) is
\[
  0\to\varprojlim_M{}^1 H^0(Q^\bullet_{w,M})
   \to H^1(\mathbf R\!\varprojlim_M Q^\bullet_{w,M})
   \to\varprojlim_M H^1(Q^\bullet_{w,M})
   \to 0.
\]
The window obstructions \([\mathfrak o^{\mathrm{ns}}_{n,M}]\)
form a class in
\(\varprojlim_M H^1(Q^\bullet_{w,M})\); they are killed by an
inverse system of windowwise counterterms
\(\{C_{n,M}\}\) precisely when this class is zero.  The Milnor
compatibility class \(\lambda_n\in
\varprojlim^1 H^0(Q^\bullet_{w,M})\) is the obstruction class for
compatibility of the windowwise counterterms across windows; it
vanishes when the truncations \(\pi_{M,N}\) are compatible with
the chosen primitives.  The combined vanishing
\(\mathfrak D^{\mathrm{ns}}_n=0\) is the fixed-order inverse-limit
criterion; the all-order statement requires the same condition for all
\(n\).
For the prototype \(\theta_3\) row, after a scalar-zero local
counterterm \(C_{\theta_3,M}\) with
\(d_M C_{\theta_3,M}=-\mathsf E_{\theta_3,M}\) is adjoined to the
finite-window complex, the coefficient \(c=1\) is windowwise primitive
and Milnor-compatible, so
\(\mathfrak D^{\mathrm{ns}}_3=0\) for that given row
(Theorem~\ref{thm:app-theta-three-finite-window-H1}).
\end{proof}

\subsection{Costello brane-defect graph QME}
\label{ssec:completion-graph-qme}

The brane-defect Costello graph QME requires a heat-kernel
propagator with controlled wavefront, removable divergences, the
RG equation, and the QME after counterterms.

\begin{defn}[Brane-defect base heat kernel and BV propagator datum]
\label{def:completion-brane-defect-propagator}
Fix a flat metric on \(\R^2\), a Hermitian metric on \(\C^2\), and
the gauge-fixing operator
\(Q^{\mathrm{GF}}=d^*_{\R^2}+\bar\partial^*_{\C^2}\) on
\(\mathcal E_H\).  The base generalized Laplacian is
\[
  D_{\mathrm{base}}
  =[d_{\R^2}+\bar\partial_{\C^2},Q^{\mathrm{GF}}]
  =\Delta_{\R^2}+\Delta_{\C^2}.
\]
Let \(K^{\mathrm{base}}_t=e^{-tD_{\mathrm{base}}}\).  The base heat
homotopy at scale \((\varepsilon,L)\) is
\[
  P^{\mathrm{base}}_{\varepsilon,L}=\int_\varepsilon^L
   \tfrac12(Q^{\mathrm{GF}}\otimes 1+1\otimes Q^{\mathrm{GF}})\,
   K^{\mathrm{base}}_t\,dt .
\]
A Costello brane-defect BV propagator is additional datum:
a coefficient coevaluation
\[
  C_{\mathrm{coeff},w}\colon
  \C\longrightarrow
  \mathcal E^{\mathrm{coeff}}_w
  \widehat\otimes
  \mathcal E^{\mathrm{coeff},\vee}_{w,\mathrm{decl}}
\]
in the declared completed coefficient category of
Definition~\ref{def:completion-host-category}, together with a clean
restriction of
\[
  P^{\mathrm{BV}}_{\varepsilon,L}
  =
  P^{\mathrm{base}}_{\varepsilon,L}\widehat\otimes
  C_{\mathrm{coeff},w}
\]
to the brane-defect support stratum \(L\times L\subset X\times X\).
Below \(P_{\varepsilon,L}\) denotes this supplied BV propagator datum.
Without \(C_{\mathrm{coeff},w}\), the displayed heat kernel is only a
base operator homotopy and does not define the BV Laplacian or Costello
graph weights.
\end{defn}

\begin{lem}[Base wavefront control and coefficient-kernel condition]
\label{thm:completion-wavefront-control}
The base heat-homotopy restriction
\((P^{\mathrm{base}}_{\varepsilon,L})|_{L\times L}\) is a
distributional kernel with wavefront set contained in
\[
  \mathrm{WF}\bigl((P^{\mathrm{base}}_{\varepsilon,L})|_{L\times L}\bigr)
  \subset
  \{((t_1,t_2),(\xi,-\xi))\,:\,t_1=t_2\},
\]
the conormal of the diagonal in \(L\times L\).  If the coefficient
coevaluation \(C_{\mathrm{coeff},w}\) of
Definition~\ref{def:completion-brane-defect-propagator} exists in the
declared completed tensor product and carries no additional base
wavefront, then \(P^{\mathrm{BV}}_{\varepsilon,L}|_{L\times L}\) has
the same base wavefront bound.  Its pairings with current insertions in
\(\mathcal D'_c(L)\widehat\otimes\mathcal P^\Pi_q\) are exactly the
C4-admissible operations: smooth test functions may multiply currents,
while products of two unregularized defect distributions are excluded.
The lemma constructs neither \(C_{\mathrm{coeff},w}\) nor the local
counterterms.
\end{lem}

\begin{proof}
The base heat kernel on \(X=\R^2\times\C^2\) factorizes as
\(K_t^{\R^2}\otimes K_t^{\C^2}\), with explicit Gaussian and
Mehler-type fundamental solutions.  Restriction to
\(L\times L=\R_t\times\R_t\) integrates
out the \(\R_s\)-direction and the holomorphic \(\C^2\)-directions,
producing a kernel on \(\R_t\times\R_t\) whose singularity is
concentrated at the diagonal.  The wavefront set is the conormal
of the diagonal by general elliptic-operator theory
\cite{hormander-vol1}.  Tensoring with a coefficient coevaluation
already convergent in the completed coefficient category changes only
the coefficient factor, not the base cotangent support.  Compatibility
with distributional current insertions is exactly C4.
\end{proof}

\begin{thm}[Graph-weight/QME criterion after supplied counterterms]
\label{thm:completion-graph-qme}
For the brane-defect interaction
\(I=I_{\mathrm{Ham}}\) of the local Hamiltonian BF theory and the
propagator \(P_{\varepsilon,L}\) of
Definition~\ref{def:completion-brane-defect-propagator}, fix the
following analytic and homological data:
\begin{enumerate}[label=\textup{(G\arabic*)}]
\item a brane-defect Costello locality theorem for this support stratum,
  giving local asymptotic expansions of the graph weights
  \(W_\Gamma(P_{\varepsilon,L},I)\) as \(\varepsilon\to0\);
\item clean wavefront restriction for the brane-defect propagator and
  current insertions, as in
  Lemma~\ref{thm:completion-wavefront-control};
\item weighted-Tate summability of the finite-scale graph expansion and
  of the local counterterm series in the chosen coefficient topology;
\item a fixed chain scalar projection, finite row arrays, and
  \(\mathfrak D^{\mathrm{ns}}_n=0\) for every order \(n\), with
  compatible local primitives \(C_{n,w,I}\).
\end{enumerate}
Then the supplied counterterm system
\(C_{\hbar_{\mathrm{QME}},w,I}=\sum_{n\ge1}\hbar_{\mathrm{QME}}^nC_{n,w,I}\) removes the local
divergences and the renormalized finite-scale expression
\[
  I[L]=\lim_{\varepsilon\to0}\sum_{\Gamma}
  \frac{\hbar_{\mathrm{QME}}^{b_1(\Gamma)}}{|\!\Aut\Gamma|}
  W_\Gamma(P_{\varepsilon,L},I+C_{\hbar_{\mathrm{QME}},w,I})
\]
exists in the stated weighted Tate topology.  In the corrected residual
complex it satisfies the finite-scale quantum master equation
\[
  Q_HI[L]+\tfrac12\{I[L],I[L]\}_L+\hbar_{\mathrm{QME}}\Delta_LI[L]=0 .
\]
Before the compatible primitives are inserted, the curvature is the
residual non-scalar QME class
\(\mathrm{Ob}^{\mathrm{red}}_{w,\partial,\hbar_{\mathrm{QME}}}\) of
Definition~\ref{def:completion-residual-nonscalar}.  The data
\textup{(G1)--(G4)} and the coefficient coevaluation of
Definition~\ref{def:completion-brane-defect-propagator} are the
separate inputs in the exact criterion for a QME-compatible
brane-defect kernel.  The bulk-to-defect
kernel
\(\kappa_{\hbar_{\mathrm{QME}},w,I}:C^\bullet_{\mathrm{CE,cont}}(\mathfrak g^{\mathrm{cur}}_{w,\hbar_{\mathrm{QME}},I})
\to\mathfrak Q^\bullet_{w,\partial,\hbar_{\mathrm{QME}}}(I)\) is defined on
generators by \(u_{a(t)dt\otimes f}\mapsto B^{\mathrm{BV}}_f(a)\),
\(c_{B,\rho}\mapsto\Theta^{\mathrm{BV}}_\rho(B)\), and the
counterterm \(C_{\hbar_{\mathrm{QME}},w,I}=\sum_n\hbar_{\mathrm{QME}}^n C_{n,w,I}\) of
Theorem~\ref{thm:completion-milnor-obstruction-vector} produces the
QME-compatible kernel
\(K_{\hbar_{\mathrm{QME}},w,I}=
\kappa_{\hbar_{\mathrm{QME}},w,I}+C_{\hbar_{\mathrm{QME}},w,I}\) satisfying
\(d_{\mathrm{QME}}K_{\hbar_{\mathrm{QME}},w,I}
+\tfrac12[K_{\hbar_{\mathrm{QME}},w,I},
K_{\hbar_{\mathrm{QME}},w,I}]_{\hbar_{\mathrm{QME}}}=0\).
\end{thm}

\begin{proof}
Condition \textup{(G1)} supplies the brane-defect version of Costello's
finite-scale graph expansion
\cite{costello-renormalization}*{Thms.~A--C}; condition \textup{(G2)}
is the clean-restriction input for the boundary current labels; and
condition \textup{(G3)} is precisely the convergence statement in the
weighted coefficient topology.  Thus the renormalized graph expression
is a well-defined element of the residual local-functional complex.
Condition \textup{(G4)} is the Milnor counterterm criterion of
Theorem~\ref{thm:completion-milnor-obstruction-vector}: the windowwise
non-scalar classes and the \(\varprojlim^1H^0\) primitive-compatibility
class vanish at every order, so the compatible family
\(C_{\hbar_{\mathrm{QME}},w,I}\) exists.  Costello's RG/QME identity
\cite{costello-renormalization}*{Lem.~10.2.2, Defn.~10.3.1}, applied after these analytic
inputs and counterterms have been fixed, gives the displayed finite-scale
QME.  Before \textup{(G4)}, the same curvature is represented by
\(\mathrm{Ob}^{\mathrm{red}}_{w,\partial,\hbar_{\mathrm{QME}}}\).  Adding the
counterterm to the generator-level kernel gives
\(K_{\hbar_{\mathrm{QME}},w,I}=
\kappa_{\hbar_{\mathrm{QME}},w,I}+C_{\hbar_{\mathrm{QME}},w,I}\), and the QME identity
is exactly
\(d_{\mathrm{QME}}K_{\hbar_{\mathrm{QME}},w,I}
+\tfrac12[K_{\hbar_{\mathrm{QME}},w,I},
K_{\hbar_{\mathrm{QME}},w,I}]_{\hbar_{\mathrm{QME}}}=0\).
\end{proof}

\begin{rmk}[\(\theta_3\) prototype]
\label{rmk:completion-theta-three-prototype}
The order-three CE-source row \(\theta_3\) is the prototype: in
each finite window \(M\ni\theta_3\) the bare \(H^1=\C\cdot
\mathsf E_{\theta_3,M}\) is killed in the enlarged complex after a
scalar-zero local counterterm \(C_{\theta_3,M}\) with
\(d_M C_{\theta_3,M}=-\mathsf E_{\theta_3,M}\) is adjoined and taken
with coefficient \(c=1\); the counterterms are Milnor-compatible
because the truncations preserve the scalar normalization, and the
corresponding order-three component in a compatible
\(C_{\hbar_{\mathrm{QME}},w,I}\) is the windowwise sum
of \(C_{\theta_3,M}\) (Theorem~\ref{thm:app-theta-three-finite-window-H1}).
This constructs only the order-three finite-window component of the
criterion.  It is not the brane-defect renormalization theorem and not
an all-order counterterm tower beyond the Milnor-compatible family of
Theorem~\ref{thm:completion-milnor-obstruction-vector}.
\end{rmk}

\subsection{\texorpdfstring{Stratified factorization algebra on \((X,L)\)}{Stratified factorization algebra on (X,L)}}
\label{ssec:completion-stratified-factorization}

The closed-open datum is the stratified factorization algebra
\(\mathrm{Obs}^{\mathrm{str}}(X,L)\); Weiss descent relates the
prefactorization datum (Appendix~\ref{sec:app-factorization-current-conventions})
to the full factorization theory.

\begin{thm}[Stratified factorization on \((X,L)\)]
\label{thm:completion-stratified-factorization}
The bulk observables \(\mathrm{Obs}_X\), the brane-line observables
\(\mathrm{Obs}_L=\mathrm{Obs}_{\partial,N}\), the bulk-to-brane
collision module
\(\mathcal M_{X/L}\), the extension-by-zero functor
\(j_*:\mathrm{Obs}_{X\setminus L}\to\mathrm{Obs}_X\), the
restriction \(i^!:\mathrm{Obs}_X\to\mathcal M_{X/L}\), and the
stratified factorization product
\(\mu^{\mathrm{str}}_{U_1,\ldots,U_k;I}:\mathrm{Obs}_X(U_1)\otimes\cdots
\otimes\mathrm{Obs}_X(U_k)\otimes\mathrm{Obs}_L(I)\to\mathrm{Obs}_L(J)\)
(whenever the \(U_i\) approach the brane interval \(I\subset J\))
together form a stratified prefactorization algebra on \((X,L)\).
Weiss descent is the additional assertion that the homotopy-colimit map
\[
  \operatornamewithlimits{hocolim}_{\mathcal U\in\mathrm{Weiss}(O)}
   \mathrm{Obs}^{\mathrm{str}}(\mathcal U)
  \xrightarrow{\sim}\mathrm{Obs}^{\mathrm{str}}(O)
\]
is an equivalence for every admissible open \(O\subset X\).  After
quantization this remains governed by the graph-QME criterion data of
Theorem~\ref{thm:completion-graph-qme} and compatible descent
homotopies.
\end{thm}

\begin{proof}
The bulk factorization algebra \(\mathrm{Obs}_X\) is constructed in the
Costello--Gwilliam factorization formalism
\cite{costello-gwilliam}*{Vol.~I Secs.~6.1, 7.2} from the Hamiltonian
BF observables on \(X=\R^2\times\C^2\); the holomorphic factor satisfies
Dolbeault descent on \(\C^2\) and the topological factor satisfies
de Rham descent on \(\R^2\), so the product satisfies Weiss
descent on the mixed open
\(O=U_{\mathrm{top}}\times B_{\mathrm{hol}}\subset X\).  The
brane-line theory is the \(E_1\)-factorization algebra
\(\mathrm{Obs}_{\partial,N}\) of
Definition~\ref{def:E1-brane-algebra}, satisfying \(E_1\)-Weiss
descent on \(L=\R_t\).  The collision module
\(\mathcal M_{X/L}\) is constructed by extension by zero from
\(X\setminus L\) and restriction to \(L\) via the costalk \(i^!\);
the stratified factorization product is the bulk-to-brane
diagonal singular product in the local Hamiltonian BF theory,
extracted from Wick contractions only after the graph-QME criterion data
of Theorem~\ref{thm:completion-graph-qme} have been supplied.  These data give
the stratified prefactorization maps.  Weiss descent requires the
homotopy-colimit descent equivalence on each stratum and compatibility
of the structure maps with the \(j_*\)/\(i^!\)-adjunction; after
quantization the same statement requires the graph-QME criterion data
and descent homotopies.
\end{proof}

\subsection{\texorpdfstring{Physical \(N\to\infty\) trace states and connected cumulants}{Physical N-to-infinity trace states and connected cumulants}}
\label{ssec:completion-large-N}

The stable primitive trace algebra is promoted to a physical
\(N\to\infty\) string state only after an actual state and connected
cumulants are supplied.

\begin{defn}[Physical \(N\to\infty\) trace-state datum]
\label{def:completion-large-N-state}
A physical \(N\to\infty\) trace-state datum on a fixed factorization window is a
sequence of continuous \(\C\)-linear functionals
\[
  \omega_{N,\Omega}:\mathcal A^q_{\partial,N}(I)\longrightarrow
  \C[[\hbar_{\mathrm{QME}},\hbar_{\mathrm{W}},\epsilon]],
  \qquad N\ge1,
\]
together with scaling exponents \(\alpha\) and \(\beta(k)\).  For each
\(N\), \(\omega_{N,\Omega}\) satisfies:
\begin{enumerate}
\item \(\omega_{N,\Omega}((Q+\hbar_{\mathrm{QME}}\Delta+\{I,-\})a)=0\)
  (Ward identity);
\item \(\omega_{N,\Omega}(ab)=(-1)^{|a||b|}\omega_{N,\Omega}(ba)\)
  (cyclicity, up to a declared anomaly line);
\item
  \(\omega_{N,\Omega}(\mu_{I_1,\ldots,I_k}(a_1\otimes\cdots\otimes a_k))
   =\omega_{N,\Omega}(a_1\cdots a_k)\)
  (factorization compatibility).
\end{enumerate}
The connected cumulants of the sequence are
\(\kappa^N_k(a_1,\ldots,a_k)=\sum_{\pi\in\mathrm{Part}(k)}
(-1)^{|\pi|-1}(|\pi|-1)!\prod_{B\in\pi}\omega_{N,\Omega}
(\prod_{i\in B}a_i)\).  The datum has an \(N\to\infty\) cumulant limit on insertions
\(f_1,\ldots,f_k\) when
\(\kappa^\infty_k(f_1,\ldots,f_k)=\lim_{N\to\infty}N^{\beta(k)}
\kappa^N_k(J_N(f_1),\ldots,J_N(f_k))\) for some scaling exponent
\(\beta(k)\) and rescaling \(J_N(f)=N^{-\alpha}\Tr f(\phi_1,\phi_2)\)
exists.
\end{defn}

\begin{thm}[Physical \(N\to\infty\) trace-state data are additional structure]
\label{thm:completion-large-N}
The assembled datum \(\mathfrak U_{\mathcal T}\) does not by itself
construct the trace-state datum of
Definition~\ref{def:completion-large-N-state} or the limits
\(\kappa^\infty_k\).  A physical \(N\to\infty\) sector consists of that
additional datum and convergence of the rescaled connected cumulants.
\end{thm}

\begin{proof}
The cyclic trace, the QME counterterms, and the weighted Köthe
estimate are necessary algebraic constraints on a possible state.
They are not a measure, a Hamiltonian, a localization cycle,
an asymptotic expansion, or a tightness theorem.  These analytic
structures belong to the physical \(N\to\infty\) trace-state theory.
\end{proof}

\subsection{The comparison datum and master Maurer--Cartan equation}
\label{ssec:completion-universal-datum}

\begin{defn}[Algebraic comparison datum and state extension]
\label{def:completion-universal-datum}
The algebraic comparison datum is
\[
  \mathfrak U_{\mathcal T}
  =(\mathsf C,\mathcal E_{\mathcal T},\mathcal E_H,
   \mathrm{SDR}_{\mathcal T},\mathcal F^q_{\mathrm{bulk}},
   \mathcal A^q_{\partial,N},\mathcal D_{\mathrm{form}},
   \Phi^q_N,K_{\hbar_{\mathrm{QME}},w,I},
   \mathcal R_{\mathrm{rad}},\mathcal P^\Pi_q,\mathcal C_{\mathcal T}),
\]
where \(\mathsf C\) is the filtered nuclear Tate presentation of the
fixed admissible host datum
(\S\ref{ssec:completion-host-category} and
Definition~\ref{defn:admissible-host-category}); \((\mathcal E_{\mathcal T},
\mathcal E_H)\) are the parent and local Hamiltonian BV complexes
(\S\ref{ssec:completion-bv-theorem},
\S\ref{ssec:completion-protected-sector-sdr});
\(\mathrm{SDR}_{\mathcal T}\) is the protected-sector cyclic BV
strong deformation retract
(Definition~\ref{def:completion-protected-sector-sdr});
\(\mathcal F^q_{\mathrm{bulk}}\) is the renormalized quantum bulk
factorization datum obtained only after the graph-QME criterion data
are supplied
(Theorem~\ref{thm:completion-stratified-factorization} and
Theorem~\ref{thm:completion-graph-qme});
\(\mathcal A^q_{\partial,N}\) is the quantum open Dirac brane
factorization algebra
(Theorem~\ref{thm:completion-tau-op-mc});
\(\mathcal D_{\mathrm{form}}\) is the admissible continuous
filtered twisted-formality, Poincar\'e--Birkhoff--Witt, and Roos
datum of Definition~\ref{def:completion-centre-cochains};
\(\Phi^q_N\) is a quantum closed-to-open derived-centre
Maurer--Cartan lift extending the fixed-rank boundary cochains
(\S\ref{ssec:completion-hochschild-centrality});
\(K_{\hbar_{\mathrm{QME}},w,I}\)
is the candidate brane-defect QME kernel obtained from the supplied
counterterm tower
(Theorem~\ref{thm:completion-graph-qme});
\(\mathcal R_{\mathrm{rad}}\) is the conditional all-bidegree
Weyl/radial normal-form criterion
(Theorem~\ref{thm:completion-omega-rad-vanishes});
\(\mathcal P^\Pi_q\) is the native pro-Matlis current target
(\S\ref{ssec:completion-promatlis}); \(\mathcal C_{\mathcal T}\)
is the matched-conventions/global-descent datum
(\S\ref{app:comparison-conventions}).  A physical \(N\to\infty\)
trace-state extension is
the enlarged datum
\[
  \mathfrak U^\omega_{\mathcal T}
  =
  \bigl(\mathfrak U_{\mathcal T},
  \{\omega_{N,\Omega}\}_{N\ge1},
  \{\alpha,\beta(k)\}_{k\ge1},
  \{\kappa^\infty_k\}_{k\ge1}\bigr),
\]
when the state sequence, scaling exponents, and connected cumulant
limits of Definition~\ref{def:completion-large-N-state} have been
supplied.  The algebraic datum \(\mathfrak U_{\mathcal T}\) does not
construct this extension.
\end{defn}

\begin{defn}[Assembled deformation complex]
\label{def:completion-universal-deformation-complex}
The algebraic assembled target deformation complex is the product filtered
\(L_\infty\) algebra
\[
  \mathfrak K^{\mathrm{alg}}_{\mathrm{univ},\mathcal T}
  =\mathfrak K_{\mathrm{prot},\mathcal T}
   \times\mathfrak K_{\mathrm{form}}
   \times\mathfrak K_{\mathrm{cent}}
   \times\mathfrak K_{\mathrm{bar/cobar}}
   \times\mathfrak K_{\mathrm{QME}}
   \times\mathfrak K_{\mathrm{rad}}
   \times\mathfrak K_{\mathrm{BM}}
   \times\mathfrak K_{\mathrm{glob},\mathcal T},
\]
with
\[
\begin{aligned}
  \mathfrak K_{\mathrm{prot},\mathcal T}
   &=\mathrm{Def}_{\mathrm{cycBV}}
     (\mathcal E_H\rightleftarrows\mathcal E_{\mathcal T}),\\
  \mathfrak K_{\mathrm{form}}
   &=\mathcal C^\bullet_{\mathcal U}
     \times\mathrm{Def}_{\mathrm{PBW}}
       (A_\hbar,\widehat U_{\hbar_{\mathrm{Rees}}}(\mathfrak h_q))
     \times
       \mathfrak C^\bullet_{\mathrm{Roos}}
       (\mathcal D_{\mathrm{form}}),\\
  \mathfrak K_{\mathrm{cent}}
   &=\mathrm{Conv}^{\mathrm{cts}}\bigl(
     C^{\mathrm{CE}}_\bullet(\mathfrak g_I),
     CC^\bullet(\mathrm{Obs}_{\partial,N},\mathrm{Obs}_{\partial,N})\bigr),\\
  \mathfrak K_{\mathrm{bar/cobar}}
   &=\mathrm{Cone}\bigl(
     \Omega C^{\mathrm{CE}}_\bullet(\mathfrak g_{\mathrm{adm}})
     \to U(\mathfrak g_{\mathrm{adm}})\bigr)\\
   &\qquad{}\oplus
     \mathrm{Def}(A_{\mathrm{op}},C_{\mathrm{op}},\tau_{\mathrm{op}}),\\
  \mathfrak K_{\mathrm{QME}}
   &=\mathfrak Q^\bullet_{w,\partial,\hbar_{\mathrm{QME}}}(I),\\
  \mathfrak K_{\mathrm{rad}}
   &=\prod_{a,b}\mathrm{Cone}(B_{a,b})[-1],\\
  \mathfrak K_{\mathrm{BM}}
   &=\mathbf R\!\varprojlim_N CC^\bullet(P_{\le N},P_{\le N}),\\
  \mathfrak K_{\mathrm{glob},\mathcal T}
   &=\mathbf R\Gamma(Y,\mathcal K_{\mathcal T}).
\end{aligned}
\]
The formality component has curvature vector
\[
  \operatorname{Ob}_{\mathrm{form}}
  =
  \bigl(
    \{\operatorname{Ob}_{\mathcal U,m+1}\}_{m\geq0},
    \operatorname{Ob}_{\mathrm{PBW}},
    \operatorname{Ob}_{\mathrm{Roos,form}}
  \bigr)
  \in H^1(\mathfrak K_{\mathrm{form}}).
\]
Its three entries are respectively the filtered
\(L_\infty\)-morphism curvature, the PBW equivalence defect, and the
transition-homotopy defect.
For the state-extended datum \(\mathfrak U^\omega_{\mathcal T}\) the
deformation complex is
\[
  \mathfrak K^{\omega}_{\mathrm{univ},\mathcal T}
  =
  \mathfrak K^{\mathrm{alg}}_{\mathrm{univ},\mathcal T}
  \times\mathfrak K_{N\to\infty},
  \qquad
  \mathfrak K_{N\to\infty}
   =
   \mathrm{Def}\bigl(\{\omega_{N,\Omega}\}_{N\ge 1},
     \{\kappa^\infty_k\}_{k\ge1}\bigr).
\]
\end{defn}

\begin{thm}[Closed-open formal Hamiltonian brane table]
\label{thm:completion-universal-formal-brane}
Let \(\mathcal T\) be an admissible protected mixed holomorphic-topological
BV/factorization target with assembled datum \(\mathfrak U_{\mathcal T}\)
of Definition~\ref{def:completion-universal-datum}.  The local
theorem of Theorem~\ref{thm:formal-local-universal-property}
produces a transported local trace-sector datum
\(\Theta_{\mathrm{loc}}\).  The realization of this datum in
\(\mathcal T\) is governed by the following component obstruction
table:
\begin{enumerate}
\item Strict algebraic realization (\(\tau=0\)): all algebraic component
  obstructions vanish and the comparison-cone contraction exists.
  Then the formal Darboux brane stalk of \(\mathcal T\) realizes the
  local Hamiltonian BF/Dirac-brane trace-sector model on
  \(\R^2_{\mathrm{top}}\times\C^2_{0,\mathrm{hol}}\).  If the
  state-extended datum \(\mathfrak U^\omega_{\mathcal T}\) is imposed,
  strict physical realization also requires
  \(\operatorname{Ob}_{N\to\infty}=0\).
\item Algebraic counterterm realization: \(\tau=\sum\tau_i\) with each
  \(\tau_i\in F^1\mathfrak K^0_i\) cancelling the \(i\)-th algebraic
  component curvature; concrete constructions are
  \(\tau_{\mathrm{rad}}\) (decorated PBW-Stokes, by
  Theorem~\ref{thm:completion-omega-rad-vanishes}),
  \(\tau_{\mathrm{QME}}=C^{\mathrm{red}}_{\hbar_{\mathrm{QME}},w}\) (Milnor
  counterterm, by
  Theorem~\ref{thm:completion-milnor-obstruction-vector}),
  \(\tau_{\mathrm{cent}}\) (Hochschild equivariance homotopies, by
  Theorem~\ref{thm:completion-centrality-homotopies}).
\item Projective/descent-torsor realization: the obstruction lands in a
  central summand and the corrected element solves the master
  equation in a central extension; the scalar Capelli anomaly
  \(\hbar_{\mathrm{W}}N[\bar c]\) is of this type.  The global Hamiltonian period
  \(\mathrm{per}_X(v)\) is a descent torsor unless a degree-two
  central extension is constructed.
\item Target replacement: the curvature is non-absorptive and
  requires a substitution; the attempted strict all-window
  Bochner--Martinelli current transfer is of this type, with
  replacement \(\mathcal P\rightsquigarrow\mathcal P^\Pi_q\) by
  Theorem~\ref{thm:completion-native-promatlis}.
\end{enumerate}
The fixed-rank trace datum is
\(J_N^s(f)=\overline{\Tr s(f)(\phi_1,\phi_2)}\) on \(\bar A\),
with finite-\(N\) trace relations and scalar contact.  Effective
scalar reduction and degreewise stable passage give
\[
  J^{\mathrm{st}}\colon\mathfrak h_q\longrightarrow
  A_{\partial,H}^{\mathrm{st}},
  \qquad
  J^{\mathrm{st}}(f)
  =\bigl(J^{\mathrm{st}}_{\leq W}
    (\operatorname{pr}_{\leq W}f)\bigr)_W,
\]
where \(\operatorname{pr}_{\leq W}\) is a Taylor coefficient
projection.  Theorem~\ref{thm:completion-shear-determination}
selects this stable trace map.  The cotangent module is the Matlis
principal-part module \(\mathcal P^\Pi_q\).  The quantum
closed-to-open derived-centre map is defined from the supplied
\(\mathcal D_{\mathrm{form}}\), centrality, and quantum master
equation data:
\(\Phi^q_N:\mathrm{Obs}^q_{\mathrm{closed},H}|_L\to
Z^q_{E_1}(\mathrm{Obs}^q_{\mathrm{open},N})\) of
\S\ref{ssec:completion-hochschild-centrality}.
\end{thm}

\begin{proof}
Each algebraic component complex of
Definition~\ref{def:completion-universal-deformation-complex} is
the deformation complex constructed in the corresponding section,
and the local datum \(\Theta_{\mathrm{loc}}\) is the assembled image
of the stable trace map \(J^{\mathrm{st}}\), the fixed-rank trace
classes \(J_N^s\), the closed Hamiltonian sector
\(\mathfrak g\), the brane factorization algebra
\(\mathcal A^{\mathrm{cl}}_{\partial,N}\), and the protected SDR
\(i_{\mathcal T}\).  The algebraic obstruction classes are
\[
\bigl(\mathrm{Ob}_{\mathrm{prot}},\mathrm{Ob}_{\mathrm{form}},
\mathrm{Ob}_{\mathrm{cent/QME/desc}},
\Omega^{\mathrm{rad}},\mathrm{Ob}^{\mathrm{red}}_{w,\partial,\hbar_{\mathrm{QME}}},
\mathrm{Ob}^\Pi_{\mathrm{BM}},\mathrm{per}_X(v),
\hbar_{\mathrm{W}}N[\bar c]\bigr).
\]
The physical state extension adds
\(\mathrm{Ob}_{N\to\infty}\).  A single curvature class exists only
after the appropriate common deformation complex and comparison maps
have been constructed; for \(\mathrm{Ob}_{N\to\infty}\) this means the
state-extended complex
\(\mathfrak K^\omega_{\mathrm{univ},\mathcal T}\), not the algebraic
complex alone.  The realization modes are read component-by-component,
with the explicit twists given by the
constructions of \S\S\ref{ssec:completion-decorated-pbw-stokes}--\S\ref{ssec:completion-large-N}.
The stable formal trace map is
Theorem~\ref{thm:completion-shear-determination}; the pro-Matlis
cotangent module is Theorem~\ref{thm:completion-native-promatlis};
the derived-centre lift assembles the supplied formality/PBW/Roos
datum, the Hochschild equivariance homotopies
(Theorem~\ref{thm:completion-centrality-homotopies}), and the
graph-QME criterion data
(Theorem~\ref{thm:completion-graph-qme}).
\end{proof}

\subsection{Compact comparison reading}
\label{ssec:completion-cross-target}

The local Hamiltonian BF/Dirac-brane theorem
(Theorem~\ref{thm:completion-universal-formal-brane}) gives the
stable formal-Darboux coordinate map
\[
  \Phi_{\mathrm{coord}}^{\mathrm{st}}\colon
  \operatorname{Regr}_u
  C^\bullet_{\mathrm{CE},\mathrm{coord}}
    (\mathfrak h_q\ltimes\mathcal P^\Pi_q[1])
  \longrightarrow
  \bigl(\mathrm{PV}_{\mathrm{adm}}
    (A_{\partial,H}^{\mathrm{st}}),
    d_{\pi_{\mathrm{lin}}}\bigr)
\]
with
\[
  c^\rho\longmapsto\partial_\rho,\qquad
  u_f\longmapsto O_f=J^{\mathrm{st}}(f),\qquad
  d_{\mathrm{CE}}u_f\longmapsto X_f.
\]
The ordinary Hochschild--Kostant--Rosenberg map realizes the
zero-Poisson special fibre
\((\mathrm{PV}_{\mathrm{adm}},0)\).  The supplied datum
\(\mathcal D_{\mathrm{form}}\) realizes
\((\mathrm{PV}_{\mathrm{adm}}[[\hbar_{\mathrm{Rees}}]],
\hbar_{\mathrm{Rees}}d_{\pi_{\mathrm{lin}}})\) in the Hochschild
complex of the PBW deformation.  A chain-level
comparison with the brane-side derived \(E_1\)-Hochschild object is
the fixed-rank Maurer--Cartan extension of the cochains
\(M_{J_N^s(f)}\), \(\Theta^{(1)}_{f,N}\), and \(T_\rho\), governed
by the Hochschild equivariance, brane-defect quantum master equation,
and descent data of
Theorems~\ref{thm:completion-centrality-homotopies},
\ref{thm:completion-graph-qme}, and
\ref{thm:completion-stratified-factorization}.  This \(E_1\)-Hochschild
comparison is weaker than the typed mixed-HT open--closed
Deligne datum, which is the ledger containing the Swiss-cheese action
tower, Hochschild-centre chart lift, Darboux-torsor flat transport,
unreduced cotangent-current lift, anomaly/locality data, and comparison
cone of Theorem~\ref{thm:typed-global-mixed-ht-deligne};
under the binding convention of \nameref{ssec:local-theorem-conventions}
the notation \(Z^{\mathrm{der}}_{\mathrm{ch}}(A)\) is used only when
that chiral centre structure has been constructed.
When imposed as external comparison data, Volume~I constructs
\(Z^{\mathrm{der}}_{\mathrm{ch}}(A)\) algebraically through ordered
bar-cobar duality, Volume~II identifies it with the bulk observables of
a 3d holomorphic-topological boundary-bulk system, and Volume~III gives
the chiral algebra
\(A=\mathrm{SpCh}_{\Sigma_{d-1},C}\circ\Phi^{\mathrm{FA}}_d(\mathcal C)\)
	through the composite Calabi--Yau-to-chiral functor
\(\Phi_d\); the formal-Darboux trace theorem computes the formal Hamiltonian
normal-slice restriction of that bulk to a finite Chan--Paton
brane stack via the fixed-rank polynomial trace classes
\(J_N^s(f)=\overline{\Tr s(f)(\phi_1,\phi_2)}\) and the
degreewise stable trace map \(J^{\mathrm{st}}\).  Taylor windows
enter as coefficient projections of this completed map.
On that normal slice the holomorphic collision product is explicit:
Theorem~\ref{thm:formal-two-dimensional-holomorphic-singular-product}
computes the diagonal local-cohomology singular product with kernel
\(K^{(2)}_\Delta\).  The cross-target comparison below must preserve
this arity-two local cohomology operation before adding descent,
QME-counterterm data, or compact-target period data.

\begin{thm}[Cross-target formal-slice principle]
\label{thm:completion-cross-target-slice}
Let \(\mathcal C\) be a CY-\(d\) category satisfying the
hypotheses of the Volume~III functor
\(\Phi_d=\mathrm{SpCh}_{\Sigma_{d-1},C}\circ\Phi^{\mathrm{FA}}_d\),
and let \(A=\Phi_d(\mathcal C)\).  Assume the Volume~II
holomorphic-topological realization \(\mathcal T_A\) admits an
admissible protected Hamiltonian normal sector at a brane \(L\),
with formal holomorphic-symplectic normal slice
\((\widehat{\C^2}_0,dz_1\wedge dz_2)\) and a \(GL_N\)-Chan--Paton
Dirac brane formal-stalk chart.  Then the protected formal-normal
trace sector of
\(Z^{\mathrm{der}}_{\mathrm{ch}}(A)\simeq
\mathrm{Obs}_{\mathrm{bulk}}(\mathcal T_A)\) has the
stable shifted-cotangent coordinate model
\[
  \operatorname{Regr}_u
  C^\bullet_{\mathrm{CE},\mathrm{coord}}
    (\mathfrak h_q\ltimes\mathcal P^\Pi_q[1])
  \xrightarrow{\ \Phi_{\mathrm{coord}}^{\mathrm{st}}\ }
  \bigl(\mathrm{PV}_{\mathrm{adm}}
    (A_{\partial,H}^{\mathrm{st}}),
    d_{\pi_{\mathrm{lin}}}\bigr),
\]
determined by
\[
  c^\rho\longmapsto\partial_\rho,\qquad
  u_f\longmapsto O_f=J^{\mathrm{st}}(f),\qquad
  d_{\mathrm{CE}}u_f\longmapsto
  X_f=d_{\pi_{\mathrm{lin}}}O_f.
\]
Here \(|c^\rho|_{\mathrm{CE}}=1\) and
\(|u_f|_{\mathrm{CE}}=2\), while
\(|u_f|_{\mathrm{coord}}=0\).  Its ordinary Hochschild realization is
the zero-Poisson special fibre
\[
  (\mathrm{PV}_{\mathrm{adm}}
    (A_{\partial,H}^{\mathrm{st}}),0)
  \xrightarrow{\ \operatorname{HKR}^{\mathrm{cont}}_{\mathrm{adm},0}\ }
  C^\bullet_{\mathrm{Hoch},\mathrm{adm},0}
    (A_{\partial,H}^{\mathrm{st}},
     A_{\partial,H}^{\mathrm{st}}).
\]
Given the supplied datum \(\mathcal D_{\mathrm{form}}\), its
linear-Poisson Hochschild realization is
\[
  \bigl(\mathrm{PV}_{\mathrm{adm}}
    (A_{\partial,H}^{\mathrm{st}})
    [[\hbar_{\mathrm{Rees}}]],
    \hbar_{\mathrm{Rees}}d_{\pi_{\mathrm{lin}}}\bigr)
  \xrightarrow{\ \mathcal U^{\hbar\pi}\ }
  \bigl(C^\bullet_{\mathrm{Hoch},\mathrm{adm}}
    (A_\hbar,A_\hbar),b_\hbar\bigr).
\]
The fixed-rank chart supplies \(J_N^s(f)\) and the chosen Koszul
Hamiltonian derivations \(X_{f,N}\); their extension to a derived
centre is the Maurer--Cartan criterion of
Theorem~\ref{thm:completion-centrality-homotopies}.  Effective scalar
reduction uses the \(\mathfrak{pgl}_N/\mathfrak{sl}_N[1]\) Koszul
model, removes \(\operatorname{Tr}\psi\) and the empty trace, and
retains \(\hbar_{\mathrm W}N[\bar c]\) as the projective central
class governed by \(\sigma_{\mathrm{sc}}\).
The cross-target comparison beyond this formal trace sector is
governed by the native obstruction classes
\[
  \bigl(\hbar_{\mathrm{W}}N[\bar c],\,
   \operatorname{Ob}_{\mathrm{form}},\,
   \Omega^{\mathrm{rad}}_{a,b},\,
   [\operatorname{Ob}^{\mathrm{red}}_{w,\partial,\hbar_{\mathrm{QME}}}],\,
   \operatorname{Ob}_{\mathrm{cent/QME/desc}},\,
   \operatorname{Ob}^\Pi_{\mathrm{BM}},\,
   \operatorname{Ob}_{\Phi_d\text{-descent}}\bigr)
\]
of Definition~\ref{def:completion-universal-deformation-complex}.
The common filtered deformation complex and its comparison maps
assemble these entries into a curvature vector.  The typed global
Deligne ledger \eqref{eq:typed-total-obstruction-ledger} supplies the
global extension criterion.
\end{thm}

\begin{proof}
The Volume~II identification
\(Z^{\mathrm{der}}_{\mathrm{ch}}(A)\simeq
\mathrm{Obs}_{\mathrm{bulk}}(\mathcal T_A)\) gives the bulk
factorization algebra; the protected-sector SDR
(\S\ref{ssec:completion-protected-sector-sdr}) restricts this
bulk to the formal Hamiltonian normal slice \(\mathcal E_H\); the
shifted-cotangent CE/PV model
(Theorem~\ref{thm:formal-disk-completed-ce-pv}) computes the
algebraic trace sector of \(\mathcal E_H\).  Theorem~
\ref{thm:completion-shear-determination} identifies
\(u_f\mapsto O_f=J^{\mathrm{st}}(f)\); the coordinate cotangent
identity gives \(c^\rho\mapsto\partial_\rho\); and the cochain
differential gives
\(d_{\mathrm{CE}}u_f\mapsto
X_f=d_{\pi_{\mathrm{lin}}}O_f\).
Proposition~\ref{prop:completion-hkr-zero-poisson} gives the ordinary
special fibre, while \(\mathcal D_{\mathrm{form}}\) gives the
linear-Poisson Hochschild comparison.  Effective scalar reduction and
\(\sigma_{\mathrm{sc}}\) place the surviving projective class in its
central summand.  The obstruction vector is the
component table of
Theorem~\ref{thm:completion-universal-formal-brane} augmented by
the \(\Phi_d\)-descent component recording the matched-conventions
obstruction (\S\ref{app:comparison-conventions}).  This proves the
formal-slice comparison statement only; a global mixed-HT Deligne
quasi-isomorphism requires the typed obstruction datum and cone
contraction of Theorem~\ref{thm:typed-global-mixed-ht-deligne}.
\end{proof}

\begin{rmk}[Modular and Borcherds/BKM bridges]
\label{rmk:completion-modular-bkm}
The local Hamiltonian Weyl scalar cocycle \(\hbar_{\mathrm{W}}N[\bar c]\) and the
Borcherds weight \(\kappa_{\mathrm{BKM}}\) of the level-two modular
section \(\Delta_5\) are distinct data.  Their identification is an
operator-level trace statement; scalar weights are protected trace
invariants of such a cocycle.  The Borcherds/BKM comparison in this formal
neighbourhood is the matched-conventions comparison of scalar invariants.
Likewise, the BCOV holomorphic anomaly, Gopakumar--Vafa,
Donaldson--Thomas, and MNOP classes are compact-target comparison data
governed, when they invoke the local stalk, by
Theorem~\ref{thm:completion-cross-target-slice} and
\S\ref{app:comparison-conventions}.  Globalization is a target-specific
obstruction problem: descent supplies only the matched-conventions
comparison datum, while a mixed-HT open--closed quasi-isomorphism further
requires the Swiss-cheese action, Morita centre lift,
Hochschild-centre lift, Darboux-torsor flat transport, unreduced
cotangent-current lift, brane-link descent, QME counterterm tower,
anomaly/locality, and comparison-cone data.
\end{rmk}

\subsection{Obstruction classes}
\label{ssec:completion-summary}

\begin{thm}[Obstruction classes of the assembled datum]
\label{thm:completion-curvature-decomposition}
The algebraic datum \(\mathfrak U_{\mathcal T}\) of
Definition~\ref{def:completion-universal-datum} has the first seven
named obstruction components in the table below.  The physical
state-extended datum \(\mathfrak U^\omega_{\mathcal T}\) adds the
eighth component \(\operatorname{Ob}_{N\to\infty}\).  They become one
curvature vector only inside the corresponding common deformation
complex with comparison maps.
\begin{center}
\small
\renewcommand{\arraystretch}{1.25}
\begin{tabularx}{\textwidth}{@{}>{\raggedright\arraybackslash}p{0.25\textwidth}
>{\raggedright\arraybackslash}X
>{\raggedright\arraybackslash}p{0.18\textwidth}
>{\raggedright\arraybackslash}p{0.23\textwidth}@{}}
Class & Construction & Mode & Reference\\\hline
\(\hbar_{\mathrm{W}}N[\bar c]\)
 & scalar Capelli cocycle on \(\bar A\)
 & projective
 & Theorem~\ref{thm:app-scalar-anomaly-transgression}, this app.~\S\ref{ssec:completion-summary}\\
\(\Omega^{\mathrm{rad}}_{a,b}\)
 & decorated PBW--Stokes contracting homotopy
 & counterterm (vanishing)
 & Theorem~\ref{thm:completion-omega-rad-vanishes}\\
\(\scriptstyle [\operatorname{Ob}^{\mathrm{red}}_{w,\partial,\hbar_{\mathrm{QME}}}]\)
 & Milnor obstruction vector \(\mathfrak D^{\mathrm{ns}}\)
 & counterterm (Milnor-compatible)
 & Theorem~\ref{thm:completion-milnor-obstruction-vector}\\
\(\operatorname{Ob}_{\mathrm{cent/QME/desc}}\)
 & centrality, brane-defect QME, descent data
 & four-component lift criterion
 & Theorem~\ref{thm:app-unreduced-cotangent-lift-obstruction}\\
\(\operatorname{Ob}^\Pi_{\mathrm{BM}}\)
 & pro-Matlis weighted K\"othe \(\mathcal P^\Pi_q\)
 & target replacement
 & Theorem~\ref{thm:completion-native-promatlis}\\
\(\operatorname{per}_X(v)\)
 & Hamiltonian descent torsor
 & descent
 & \S\ref{ssec:completion-cross-target}\\
\(\operatorname{Ob}_{\mathrm{prot}}\)
 & cyclic BV strong deformation retract
 & \(L_\infty\)-transferred bracket / counterterm
 & Theorem~\ref{thm:completion-prot-sector-obstruction}\\
\(\operatorname{Ob}_{N\to\infty}\)
 & trace-state datum; cumulant convergence
 & trace-state extension
 & Theorem~\ref{thm:completion-large-N}
\end{tabularx}
\end{center}
The hypotheses of the local theorem are expressed as explicit
obstruction classes, relative constructions, or replacement
data in Theorem~\ref{thm:completion-curvature-decomposition}.  The
remaining hypotheses are the matched-conventions data
\(\mathcal C_{\mathcal T}\); the protected strong deformation retract;
the all-bidegree radial primitive outside the verified finite windows;
the all-order QME counterterm compatibility; and the physical
\(N\to\infty\) trace-state datum.
\end{thm}

\begin{proof}
Each row is a corollary of the cited theorem, with the realization
mode read off from the curvature mechanism: counterterm-vanishing
via contracting homotopy when fixed (radial), Milnor counterterm-compatibility
(QME), the four-component centre/QME/descent lift criterion,
pro-target replacement
(BM), central extension (Capelli scalar), descent torsor (period), and SDR
construction (protected sector).  The
matched-conventions datum is controlled by
\S\ref{app:comparison-conventions}; the additional analytic state data
are those named in the theorem.
\end{proof}

\begin{rmk}[Remaining data]
\label{rmk:completion-remaining-data}
After Theorem~\ref{thm:completion-curvature-decomposition}, the remaining data are:
(i) the parent BV theory
\((\mathcal E_{\mathcal T},Q_{\mathcal T},\omega_{\mathcal T},
S_{\mathcal T})\) and the protected supercharge \(\mathcal Q\) for
a specific compact target \(\mathcal T\); (ii) the
matched-conventions datum \(\mathcal C_{\mathcal T}\) (residue
versus equivariant Euler normalization, orientation, framing on
\(S^3\), genus expansion endpoints) for that target;
(iii) the uniform decorated PBW-Stokes contracting homotopy
\((A^{(j)}_{a,b,N},\tau_{a,b})\) at every bidegree, of which the
finite-window checks give the recorded ranges
(Remark~\ref{rmk:completion-omega-rad-status}); and (iv) the
Statement~\ref{stmt:app-radial-parts-hypotheses}(1)--(2)
homomorphism and injectivity theorems from the cited radial-parts
literature; (v) the all-order QME counterterm compatibility across the
inverse system; and (vi) the physical \(N\to\infty\) trace-state datum and cumulant
convergence data.
The remaining clauses are represented by constructed data or by
obstruction classes in the algebraic assembled datum
\(\mathfrak U_{\mathcal T}\), except for the physical \(N\to\infty\) row,
which lives only after the state extension \(\mathfrak U^\omega_{\mathcal T}\)
has been supplied.
\end{rmk}

\subsection{Proofs and obstruction classes}
\label{ssec:completion-construction-status}

Each component of the assembled datum is classified by the theorem that
proves it, the hypotheses under which it is formulated, and the
cohomology class that obstructs its extension.

\begin{description}
\item[Theorem~\ref{thm:completion-cme} (classical master equation).]
  Proved from first principles in the filtered nuclear Tate presentation
  \(\mathsf C\) of \(\mathcal B^{\mathrm{adm}}\), using \(D^2=0\),
  graded Leibniz, Jacobi in
  \(\mathfrak h\), and continuity of the residue pairing.  The proof is
  internal to that presentation.

\item[Theorem~\ref{thm:completion-shear-determination} (trace
  uniqueness).]  Proved from the linear normalization
  \(c_{1,0}=c_{0,1}=1\) and Darboux-naturality.  The
  nonlinear Hamiltonian shear \(H_k=z_2^{k+1}/(k+1)\) closes the
  proof by giving the multiplicative relation.

\item[Theorem~\ref{thm:completion-omega-rad-vanishes} (decorated
  radial reformulation).]  This is a reformulation toward
  first-principles closure.  The all-bidegree statement is
  converted into a primitive-existence problem on a finite-rank
  decorated bicomplex; exact rational calculation gives the primitives
  at the recorded bidegrees, and the all-\((a,b)\) statement is the
  assertion that this primitive construction is uniform.  The remaining
  obstruction is the
  all-bidegree decorated primitive.

\item[Theorem~\ref{thm:completion-native-promatlis} (pro-Matlis
  target).]  Proved as a replacement target with continuous action
  and modified window coherence.  The continuity input is the explicit
  residue-coadjoint coefficient formula and seminorm inequality
  \(p^*_n(f\cdot\rho)\le
  C_n p_{n+r}(f)p^*_{n+r}(\rho)\) of
  Proposition~\ref{prop:completion-promatlis-continuity}.  Terminality
  in \(\mathrm{Curr}_{\mathrm{BM}}\) is not asserted.

\item[Theorem~\ref{thm:completion-p0-lift} and
  Theorem~\ref{thm:completion-mc-kernel} (CE/PV levels 2 and 3).]
  Proved from the bracket-admissible \(B_q\) of
  Definition~\ref{def:completion-bracket-admissible} and the
  kernel-admissible \(\mathcal K_{B_q}\) of
  Definition~\ref{def:completion-kernel-admissible}, both
  constructed in the K\"othe topology.  Soundness depends on the
  same seminorm estimate of
  Proposition~\ref{prop:completion-promatlis-continuity}, together with
  the finite-window Casimir convergence of
  Definition~\ref{def:completion-kernel-admissible}.

\item[Theorem~\ref{thm:completion-tate-bar-cobar} (CE/PV level 4,
  bar-cobar).]  Proved at each finite quotient by Quillen
  rational-homotopy bar-cobar.  The completed comparison requires the
  derived-limit obstruction terms to vanish and the bar-cobar functor
  to commute with the chosen completion.

\item[Theorem~\ref{thm:completion-tau-op-mc} (CE/PV level 5,
  open-trace \(A_\infty\)-Koszul).]  Proved by Koszul duality in
  the cyclic \(A_\infty\) operad (Loday--Vallette~\cite{loday-vallette},
  Getzler--Jones~\cite{getzler-jones-aoperad}).  Soundness depends on the cyclicity hypothesis
  on \(A_{\mathrm{op}}\), which holds for the brane factorization
  algebra \(A^{\mathrm{cl}}_{\partial,N}\).

\item[Theorem~\ref{thm:completion-centrality-homotopies} (Hochschild
  equivariance, bracket-admissible sector).]  On the
	  bracket-admissible sector \(B_q\), the Koszul homotopy for
	  \([\phi_1,\phi_2]=0\) gives the explicit
  Hochschild homotopies.  The unrestricted centre lift on
  arbitrary cotangent currents requires the brane-defect
  propagator extension (Theorem~\ref{thm:completion-graph-qme}),
  acknowledged in Remark~\ref{rmk:completion-cotangent-centrality}.

\item[Theorem~\ref{thm:completion-prot-sector-obstruction}
  (protected-sector SDR obstruction).]  Proved as an obstruction
  theorem given an SDR \((i_T,p_T,h_T)\); existence of the SDR
  for a specific compact target is a theorem in the compact field
  theory, constrained by the matched-conventions obstruction.

\item[Theorem~\ref{thm:completion-sigma-sc} (\(\sigma_{\mathrm{sc}}\)
  filtered chain map).]  The unreduced scalar cocycle has primitive
  \(\eta\).  On the reduced Hamiltonian algebra the class
  \([\bar c]\) remains projective.  The filtered lift is governed by
  the obstruction tower of
  Theorem~\ref{thm:completion-milnor-obstruction-vector}.

\item[Theorem~\ref{thm:completion-milnor-obstruction-vector}
  (Milnor obstruction vector).]  The Milnor exact sequence is applied
  to the inverse system of finite Hamiltonian windows.  Vanishing of
  \(\mathfrak D^{\mathrm{ns}}_n\)
  for every \(n\) is the all-order QME counterterm criterion; the
  \(\theta_3\) row verifies this at \(n=3\).

\item[Theorem~\ref{thm:completion-graph-qme} (Costello brane-defect
  graph-QME criterion).]  This is a criterion after four supplied
  inputs: the brane-defect finite-scale graph expansion, clean
  wavefront restriction, weighted-Tate summability, and vanishing of the
  Milnor counterterm obstruction tower.  Under those data the
  counterterm-corrected finite-scale interaction has vanishing QME
  curvature.  The brane-defect renormalization theorem and the weighted
  convergence estimates are part of the supplied four-input criterion;
  without the compatible Milnor counterterms the curvature is the
  residual non-scalar class
  \(\mathrm{Ob}^{\mathrm{red}}_{w,\partial,\hbar_{\mathrm{QME}}}\).

\item[Theorem~\ref{thm:completion-stratified-factorization}
  (stratified factorization on \((X,L)\)).]  Proved at the
prefactorization level by extension by zero plus brane-defect
  collision product.  The holomorphic part of this collision product is
  the native diagonal local-cohomology operation of
  Theorem~\ref{thm:formal-two-dimensional-holomorphic-singular-product}; the
  brane-defect term adds the interval support current and defect
  restriction.  Weiss descent follows from descent on each stratum
  and \(j_*\)/\(i^!\)-adjunction.  The post-quantization
  statement requires the graph-QME criterion data of
  Theorem~\ref{thm:completion-graph-qme}.

\item[Theorem~\ref{thm:completion-large-N} (physical \(N\to\infty\) trace-state
  datum).]  This theorem assumes a trace-state sequence, Ward identity,
  cyclicity, factorization, scaling exponents, and convergence of
  connected cumulants.  Köthe continuity and QME counterterms
  are necessary algebraic constraints, not an \(N\to\infty\) theorem.

\item[Theorem~\ref{thm:completion-universal-formal-brane} (closed-open
  formal Hamiltonian brane comparison).]  This assembles
  the preceding obstruction classes and relative constructions
  into the algebraic comparison datum \(\mathfrak U_{\mathcal T}\).  The
  remaining hypotheses are the matched-conventions datum, the SDR
  existence, the radial all-bidegree primitive outside verified
  windows, the all-order QME counterterm compatibility, and the physical
  \(N\to\infty\) trace-state datum.

\item[Theorem~\ref{thm:completion-cross-target-slice} (cross-target
  formal-slice principle).]  Proved as the formal restriction of
  the trace sector of the Volumes I/II derived chiral centre to the
  protected Hamiltonian normal sector, with the component obstruction classes recording
  the global, quantum, chain-level, all-window, and Calabi--Yau
  comparison obstructions.  It is formulated under the Volume~III
  categorical theorem and Volume~II HT realization.
\end{description}

\begin{rmk}[Algebraic and analytic hypotheses]
\label{rmk:completion-three-tier}
The classical master equation and the trace uniqueness theorem are
identities in the filtered nuclear Tate presentation \(\mathsf C\); they
become statements in \(\mathcal B^{\mathrm{adm}}\) after the operation
and kernel lists of the relevant theorem are attached.  The pro-Matlis
replacement target, the finite-quotient bar-cobar comparison, the
open-trace \(A_\infty\)-Koszul comparison, the Milnor obstruction
vector, and the universal obstruction vector carry their derived-limit
obstruction classes explicitly.  The CE/PV levels 2--5, the Hochschild
equivariance homotopies, the \(\sigma_{\mathrm{sc}}\) chain map, the
brane-defect graph QME, and the stratified factorization algebra require
bracket-admissibility and kernel-admissibility, a protected SDR, and
the Stokes hypotheses for the decorated PBW complex.  A compact target
also requires a parent BV theory and protected supercharge, the
matched-conventions datum, a uniform decorated PBW-Stokes contracting
homotopy in every bidegree, clauses (1)--(2) of
Statement~\ref{stmt:app-radial-parts-hypotheses}, all-order QME
counterterm compatibility, and a physical \(N\to\infty\) trace-state
datum.
\end{rmk}
